%%%%%%%%%%%%%%%%%%%%%%%%%%%%%%%%%%%%%%%%%%%%%%%%%%%%%%%%%%%%
% 11.5 NUMERICAL PARTIAL DIFFERENTIAL EQUATIONS
% The Composition of Advanced Mathematics
%%%%%%%%%%%%%%%%%%%%%%%%%%%%%%%%%%%%%%%%%%%%%%%%%%%%%%%%%%%%

\section{Numerical Partial Differential Equations}
\label{sec:numerical-pdes}

\prerequisite{
Partial differential equations, multivariable calculus, numerical analysis,
numerical linear algebra, numerical ODEs, Taylor expansions, boundary-value
problems, and basic programming concepts.
}

\learningobjectives{
By the end of this section, the reader should be able to:
\begin{itemize}
    \item explain why PDEs often require numerical approximation;
    \item distinguish elliptic, parabolic, and hyperbolic PDEs;
    \item construct spatial and temporal computational grids;
    \item approximate first and second derivatives using finite differences;
    \item derive standard centered finite-difference stencils;
    \item discretize the Poisson equation;
    \item apply Dirichlet, Neumann, and Robin boundary conditions;
    \item understand the discrete Laplacian;
    \item construct finite-difference matrix systems;
    \item recognize sparse matrix structure arising from PDE discretization;
    \item understand truncation error and consistency for PDE schemes;
    \item understand convergence of spatial discretizations;
    \item discretize the heat equation;
    \item derive the Forward-Time Centered-Space scheme;
    \item derive the Backward-Time Centered-Space scheme;
    \item derive the Crank--Nicolson method;
    \item analyze stability using Fourier or von Neumann analysis;
    \item understand the CFL condition;
    \item distinguish conditional and unconditional stability;
    \item understand numerical diffusion and numerical dispersion;
    \item discretize the advection equation;
    \item recognize upwind schemes;
    \item understand why centered advection schemes may be unstable;
    \item recognize Lax--Friedrichs and Lax--Wendroff methods;
    \item discretize the wave equation;
    \item derive a standard second-order wave-equation scheme;
    \item understand domain-of-dependence ideas;
    \item understand finite-volume methods;
    \item formulate conservation laws in integral form;
    \item understand numerical fluxes;
    \item recognize Godunov-type methods conceptually;
    \item understand finite-element methods conceptually;
    \item derive a weak formulation;
    \item understand trial and test spaces;
    \item recognize basis functions and assembly;
    \item understand mass and stiffness matrices;
    \item understand Galerkin approximation;
    \item recognize spectral methods;
    \item understand collocation and modal spectral approaches;
    \item recognize pseudospectral differentiation matrices;
    \item understand the method of lines;
    \item recognize operator splitting;
    \item understand adaptive mesh refinement conceptually;
    \item distinguish structured and unstructured meshes;
    \item recognize consistency, stability, and convergence issues in PDEs;
    \item understand discrete maximum principles conceptually;
    \item understand conservation and monotonicity properties;
    \item recognize nonlinear PDE discretizations;
    \item understand Newton linearization for nonlinear PDE systems;
    \item recognize iterative solvers and preconditioners for PDE matrices;
    \item connect numerical PDEs with fluid dynamics, heat transfer,
          elasticity, electromagnetics, finance, wave propagation, and
          scientific computing.
\end{itemize}
}

%%%%%%%%%%%%%%%%%%%%%%%%%%%%%%%%%%%%%%%%%%%%%%%%%%%%%%%%%%%%
\subsection{Why Numerical PDEs?}
\label{subsec:why-numerical-pdes}
%%%%%%%%%%%%%%%%%%%%%%%%%%%%%%%%%%%%%%%%%%%%%%%%%%%%%%%%%%%%

Partial differential equations describe systems whose unknowns depend on
multiple independent variables.

Examples include

\[
\boxed{
u_t=\alpha u_{xx},
}
\]

\[
\boxed{
u_{tt}=c^2u_{xx},
}
\]

and

\[
\boxed{
-\Delta u=f.
}
\]

Closed-form solutions exist only for special geometries, coefficients,
boundary conditions, and forcing terms.

Numerical PDE methods replace the continuous problem by a finite set of
algebraic equations that can be solved computationally.

%%%%%%%%%%%%%%%%%%%%%%%%%%%%%%%%%%%%%%%%%%%%%%%%%%%%%%%%%%%%
\subsection{Classification of Second-Order PDEs}
\label{subsec:pde-classification-numerical}
%%%%%%%%%%%%%%%%%%%%%%%%%%%%%%%%%%%%%%%%%%%%%%%%%%%%%%%%%%%%

A second-order PDE in two variables may contain terms

\[
A u_{xx}
+
B u_{xy}
+
C u_{yy}.
\]

The discriminant

\[
\boxed{
B^2-4AC
}
\]

helps classify the equation.

The classical categories are:

\[
\boxed{
B^2-4AC<0
\quad
\text{elliptic},
}
\]

\[
\boxed{
B^2-4AC=0
\quad
\text{parabolic},
}
\]

and

\[
\boxed{
B^2-4AC>0
\quad
\text{hyperbolic}.
}
\]

Different PDE classes require different numerical techniques.

%%%%%%%%%%%%%%%%%%%%%%%%%%%%%%%%%%%%%%%%%%%%%%%%%%%%%%%%%%%%
\subsection{Canonical Examples}
\label{subsec:pde-canonical-examples}
%%%%%%%%%%%%%%%%%%%%%%%%%%%%%%%%%%%%%%%%%%%%%%%%%%%%%%%%%%%%

A standard elliptic PDE is Poisson's equation:

\[
\boxed{
-\Delta u=f.
}
\]

A standard parabolic PDE is the heat equation:

\[
\boxed{
u_t=\alpha\Delta u.
}
\]

A standard hyperbolic PDE is the wave equation:

\[
\boxed{
u_{tt}=c^2\Delta u.
}
\]

Another important hyperbolic equation is linear advection:

\[
\boxed{
u_t+cu_x=0.
}
\]

%%%%%%%%%%%%%%%%%%%%%%%%%%%%%%%%%%%%%%%%%%%%%%%%%%%%%%%%%%%%
\subsection{Computational Grid}
\label{subsec:pde-computational-grid}
%%%%%%%%%%%%%%%%%%%%%%%%%%%%%%%%%%%%%%%%%%%%%%%%%%%%%%%%%%%%

For a one-dimensional spatial domain

\[
[a,b],
\]

define grid points

\[
\boxed{
x_i=a+ih,
\qquad
i=0,\ldots,N,
}
\]

where

\[
\boxed{
h=\frac{b-a}{N}.
}
\]

For time-dependent PDEs, define

\[
\boxed{
t^n=n\Delta t.
}
\]

The numerical approximation is written

\[
\boxed{
u_i^n
\approx
u(x_i,t^n).
}
\]

%%%%%%%%%%%%%%%%%%%%%%%%%%%%%%%%%%%%%%%%%%%%%%%%%%%%%%%%%%%%
\subsection{Finite Differences}
\label{subsec:pde-finite-differences}
%%%%%%%%%%%%%%%%%%%%%%%%%%%%%%%%%%%%%%%%%%%%%%%%%%%%%%%%%%%%

Finite-difference methods replace derivatives with combinations of nearby
grid values.

For example,

\[
\boxed{
u_x(x_i)
\approx
\frac{
u_{i+1}-u_{i-1}
}{
2h
}.
}
\]

Similarly,

\[
\boxed{
u_{xx}(x_i)
\approx
\frac{
u_{i+1}-2u_i+u_{i-1}
}{
h^2
}.
}
\]

These formulas are obtained from Taylor expansions.

%%%%%%%%%%%%%%%%%%%%%%%%%%%%%%%%%%%%%%%%%%%%%%%%%%%%%%%%%%%%
\subsection{Centered First Derivative}
\label{subsec:pde-centered-first-derivative}
%%%%%%%%%%%%%%%%%%%%%%%%%%%%%%%%%%%%%%%%%%%%%%%%%%%%%%%%%%%%

Expand

\[
u(x_i+h)
=
u_i
+
hu_i'
+
\frac{h^2}{2}u_i''
+
\frac{h^3}{6}u_i'''
+
O(h^4),
\]

and

\[
u(x_i-h)
=
u_i
-
hu_i'
+
\frac{h^2}{2}u_i''
-
\frac{h^3}{6}u_i'''
+
O(h^4).
\]

Subtracting gives

\[
u_{i+1}-u_{i-1}
=
2hu_i'
+
O(h^3).
\]

Therefore,

\[
\boxed{
u_x(x_i)
=
\frac{
u_{i+1}-u_{i-1}
}{
2h
}
+
O(h^2).
}
\]

%%%%%%%%%%%%%%%%%%%%%%%%%%%%%%%%%%%%%%%%%%%%%%%%%%%%%%%%%%%%
\subsection{Centered Second Derivative}
\label{subsec:pde-centered-second-derivative}
%%%%%%%%%%%%%%%%%%%%%%%%%%%%%%%%%%%%%%%%%%%%%%%%%%%%%%%%%%%%

Adding the two Taylor expansions gives

\[
u_{i+1}
-
2u_i
+
u_{i-1}
=
h^2u_i''
+
O(h^4).
\]

Thus,

\[
\boxed{
u_{xx}(x_i)
=
\frac{
u_{i+1}-2u_i+u_{i-1}
}{
h^2
}
+
O(h^2).
}
\]

This is one of the most important finite-difference formulas.

%%%%%%%%%%%%%%%%%%%%%%%%%%%%%%%%%%%%%%%%%%%%%%%%%%%%%%%%%%%%
\subsection{Stencil Notation}
\label{subsec:pde-stencils}
%%%%%%%%%%%%%%%%%%%%%%%%%%%%%%%%%%%%%%%%%%%%%%%%%%%%%%%%%%%%

A stencil describes the grid points used in a finite-difference formula.

For the centered second derivative,

\[
\boxed{
\frac{
u_{i-1}-2u_i+u_{i+1}
}{
h^2
},
}
\]

the one-dimensional stencil coefficients are

\[
\boxed{
\frac1{h^2}
\begin{bmatrix}
1 & -2 & 1
\end{bmatrix}.
}
\]

Stencil notation makes multidimensional schemes easier to visualize.

%%%%%%%%%%%%%%%%%%%%%%%%%%%%%%%%%%%%%%%%%%%%%%%%%%%%%%%%%%%%
\subsection{The One-Dimensional Poisson Equation}
\label{subsec:one-dimensional-poisson}
%%%%%%%%%%%%%%%%%%%%%%%%%%%%%%%%%%%%%%%%%%%%%%%%%%%%%%%%%%%%

Consider

\[
\boxed{
-u''(x)=f(x),
\qquad
a<x<b,
}
\]

with Dirichlet boundary conditions

\[
\boxed{
u(a)=\alpha,
\qquad
u(b)=\beta.
}
\]

At interior grid point \(x_i\),

\[
-u''(x_i)
\approx
-
\frac{
u_{i+1}-2u_i+u_{i-1}
}{
h^2
}.
\]

Thus,

\[
\boxed{
-u_{i-1}
+
2u_i
-
u_{i+1}
=
h^2f_i.
}
\]

%%%%%%%%%%%%%%%%%%%%%%%%%%%%%%%%%%%%%%%%%%%%%%%%%%%%%%%%%%%%
\subsection{Poisson Matrix System}
\label{subsec:poisson-matrix-system}
%%%%%%%%%%%%%%%%%%%%%%%%%%%%%%%%%%%%%%%%%%%%%%%%%%%%%%%%%%%%

For interior unknowns

\[
U
=
\begin{pmatrix}
u_1\\
u_2\\
\vdots\\
u_{N-1}
\end{pmatrix},
\]

the discrete problem becomes

\[
\boxed{
AU=b,
}
\]

where

\[
\boxed{
A
=
\frac1{h^2}
\begin{pmatrix}
2&-1&0&\cdots&0\\
-1&2&-1&\ddots&\vdots\\
0&-1&2&\ddots&0\\
\vdots&\ddots&\ddots&\ddots&-1\\
0&\cdots&0&-1&2
\end{pmatrix}.
}
\]

Boundary values contribute to the right-hand side.

%%%%%%%%%%%%%%%%%%%%%%%%%%%%%%%%%%%%%%%%%%%%%%%%%%%%%%%%%%%%
\subsection{Worked Example: One-Dimensional Poisson Problem}
\label{subsec:worked-poisson-1d}
%%%%%%%%%%%%%%%%%%%%%%%%%%%%%%%%%%%%%%%%%%%%%%%%%%%%%%%%%%%%

\begin{workedexamplebox}{Three Interior Grid Points}

Consider

\[
-u''(x)=1,
\qquad
0<x<1,
\]

with

\[
u(0)=0,
\qquad
u(1)=0.
\]

Use

\[
h=\frac14.
\]

The interior points are

\[
x_1=\frac14,
\qquad
x_2=\frac12,
\qquad
x_3=\frac34.
\]

The finite-difference equations are

\[
-u_{i-1}+2u_i-u_{i+1}
=
h^2
=
\frac1{16}.
\]

Thus,

\[
\begin{pmatrix}
2&-1&0\\
-1&2&-1\\
0&-1&2
\end{pmatrix}
\begin{pmatrix}
u_1\\
u_2\\
u_3
\end{pmatrix}
=
\frac1{16}
\begin{pmatrix}
1\\
1\\
1
\end{pmatrix}.
\]

By symmetry,

\[
u_1=u_3.
\]

Solving gives

\[
u_1=\frac{3}{32},
\qquad
u_2=\frac18,
\qquad
u_3=\frac{3}{32}.
\]

\begin{answerbox}
\[
\boxed{
U
=
\begin{pmatrix}
3/32\\
1/8\\
3/32
\end{pmatrix}.
}
\]
\end{answerbox}

\end{workedexamplebox}

%%%%%%%%%%%%%%%%%%%%%%%%%%%%%%%%%%%%%%%%%%%%%%%%%%%%%%%%%%%%
\subsection{Dirichlet Boundary Conditions}
\label{subsec:dirichlet-boundary-numerical}
%%%%%%%%%%%%%%%%%%%%%%%%%%%%%%%%%%%%%%%%%%%%%%%%%%%%%%%%%%%%

A Dirichlet condition specifies the solution value:

\[
\boxed{
u=g
\quad
\text{on the boundary}.
}
\]

These values are known directly and are inserted into the discrete
equations.

%%%%%%%%%%%%%%%%%%%%%%%%%%%%%%%%%%%%%%%%%%%%%%%%%%%%%%%%%%%%
\subsection{Neumann Boundary Conditions}
\label{subsec:neumann-boundary-numerical}
%%%%%%%%%%%%%%%%%%%%%%%%%%%%%%%%%%%%%%%%%%%%%%%%%%%%%%%%%%%%

A Neumann condition specifies the normal derivative:

\[
\boxed{
\frac{\partial u}{\partial n}
=
g.
}
\]

In one dimension at the left endpoint, one may approximate

\[
u_x(a)
\]

using a one-sided difference:

\[
\boxed{
\frac{
u_1-u_0
}{
h
}
\approx
g.
}
\]

Higher-order boundary approximations are often preferable.

%%%%%%%%%%%%%%%%%%%%%%%%%%%%%%%%%%%%%%%%%%%%%%%%%%%%%%%%%%%%
\subsection{Robin Boundary Conditions}
\label{subsec:robin-boundary-numerical}
%%%%%%%%%%%%%%%%%%%%%%%%%%%%%%%%%%%%%%%%%%%%%%%%%%%%%%%%%%%%

A Robin condition combines value and derivative:

\[
\boxed{
\alpha u
+
\beta
\frac{\partial u}{\partial n}
=
g.
}
\]

Finite differences or finite elements can incorporate this condition into
the boundary equations.

%%%%%%%%%%%%%%%%%%%%%%%%%%%%%%%%%%%%%%%%%%%%%%%%%%%%%%%%%%%%
\subsection{Two-Dimensional Laplacian}
\label{subsec:two-dimensional-laplacian}
%%%%%%%%%%%%%%%%%%%%%%%%%%%%%%%%%%%%%%%%%%%%%%%%%%%%%%%%%%%%

In two dimensions,

\[
\boxed{
\Delta u
=
u_{xx}+u_{yy}.
}
\]

For equal grid spacing \(h\),

\[
u_{xx}
\approx
\frac{
u_{i+1,j}
-
2u_{i,j}
+
u_{i-1,j}
}{
h^2
},
\]

and

\[
u_{yy}
\approx
\frac{
u_{i,j+1}
-
2u_{i,j}
+
u_{i,j-1}
}{
h^2
}.
\]

Therefore,

\[
\boxed{
\Delta u_{i,j}
\approx
\frac{
u_{i+1,j}
+
u_{i-1,j}
+
u_{i,j+1}
+
u_{i,j-1}
-
4u_{i,j}
}{
h^2
}.
}
\]

%%%%%%%%%%%%%%%%%%%%%%%%%%%%%%%%%%%%%%%%%%%%%%%%%%%%%%%%%%%%
\subsection{Five-Point Laplacian Stencil}
\label{subsec:five-point-laplacian}
%%%%%%%%%%%%%%%%%%%%%%%%%%%%%%%%%%%%%%%%%%%%%%%%%%%%%%%%%%%%

The standard two-dimensional stencil is

\[
\boxed{
\frac1{h^2}
\begin{array}{ccc}
&1&\\
1&-4&1\\
&1&
\end{array}.
}
\]

This is called the five-point stencil because it uses the center point and
its four nearest axis-aligned neighbors.

%%%%%%%%%%%%%%%%%%%%%%%%%%%%%%%%%%%%%%%%%%%%%%%%%%%%%%%%%%%%
\subsection{Discrete Poisson Equation in Two Dimensions}
\label{subsec:discrete-poisson-2d}
%%%%%%%%%%%%%%%%%%%%%%%%%%%%%%%%%%%%%%%%%%%%%%%%%%%%%%%%%%%%

For

\[
-\Delta u=f,
\]

the finite-difference equation is

\[
\boxed{
4u_{i,j}
-
u_{i+1,j}
-
u_{i-1,j}
-
u_{i,j+1}
-
u_{i,j-1}
=
h^2f_{i,j}.
}
\]

All interior nodes together form a large sparse linear system.

%%%%%%%%%%%%%%%%%%%%%%%%%%%%%%%%%%%%%%%%%%%%%%%%%%%%%%%%%%%%
\subsection{Sparse Structure}
\label{subsec:pde-sparse-structure}
%%%%%%%%%%%%%%%%%%%%%%%%%%%%%%%%%%%%%%%%%%%%%%%%%%%%%%%%%%%%

Each finite-difference equation involves only nearby grid points.

Therefore each row of the matrix contains only a small number of nonzero
entries.

For a two-dimensional five-point stencil, each interior row has at most
five important nonzeros.

Thus PDE discretization naturally produces sparse matrices.

%%%%%%%%%%%%%%%%%%%%%%%%%%%%%%%%%%%%%%%%%%%%%%%%%%%%%%%%%%%%
\subsection{Kronecker Product Structure}
\label{subsec:pde-kronecker-structure}
%%%%%%%%%%%%%%%%%%%%%%%%%%%%%%%%%%%%%%%%%%%%%%%%%%%%%%%%%%%%

On a rectangular grid, the two-dimensional discrete Laplacian can often be
written using Kronecker products.

If \(T\) is the one-dimensional tridiagonal Laplacian matrix, then

\[
\boxed{
A
=
I\otimes T
+
T\otimes I.
}
\]

This structure is valuable both analytically and computationally.

%%%%%%%%%%%%%%%%%%%%%%%%%%%%%%%%%%%%%%%%%%%%%%%%%%%%%%%%%%%%
\subsection{Discrete Maximum Principle Preview}
\label{subsec:discrete-maximum-principle}
%%%%%%%%%%%%%%%%%%%%%%%%%%%%%%%%%%%%%%%%%%%%%%%%%%%%%%%%%%%%

Continuous elliptic equations often satisfy maximum principles.

Well-designed discrete schemes may inherit analogous properties.

For example, under appropriate assumptions the discrete solution may not
develop artificial interior maxima or minima beyond those permitted by
the forcing and boundary data.

Such structure contributes to stability and physical reliability.

%%%%%%%%%%%%%%%%%%%%%%%%%%%%%%%%%%%%%%%%%%%%%%%%%%%%%%%%%%%%
\subsection{The Heat Equation}
\label{subsec:numerical-heat-equation}
%%%%%%%%%%%%%%%%%%%%%%%%%%%%%%%%%%%%%%%%%%%%%%%%%%%%%%%%%%%%

Consider

\[
\boxed{
u_t
=
\alpha u_{xx},
}
\]

where

\[
\alpha>0
\]

is the diffusivity.

Define

\[
\boxed{
u_i^n
\approx
u(x_i,t^n).
}
\]

A numerical method must approximate both time and space derivatives.

%%%%%%%%%%%%%%%%%%%%%%%%%%%%%%%%%%%%%%%%%%%%%%%%%%%%%%%%%%%%
\subsection{Forward-Time Centered-Space Method}
\label{subsec:ftcs-heat}
%%%%%%%%%%%%%%%%%%%%%%%%%%%%%%%%%%%%%%%%%%%%%%%%%%%%%%%%%%%%

Approximate

\[
u_t
\]

with forward Euler:

\[
\frac{
u_i^{n+1}-u_i^n
}{
\Delta t
}.
\]

Approximate

\[
u_{xx}
\]

with the centered difference:

\[
\frac{
u_{i+1}^n
-
2u_i^n
+
u_{i-1}^n
}{
h^2
}.
\]

Thus,

\[
\boxed{
\frac{
u_i^{n+1}-u_i^n
}{
\Delta t
}
=
\alpha
\frac{
u_{i+1}^n
-
2u_i^n
+
u_{i-1}^n
}{
h^2
}.
}
\]

%%%%%%%%%%%%%%%%%%%%%%%%%%%%%%%%%%%%%%%%%%%%%%%%%%%%%%%%%%%%
\subsection{Diffusion Number}
\label{subsec:diffusion-number}
%%%%%%%%%%%%%%%%%%%%%%%%%%%%%%%%%%%%%%%%%%%%%%%%%%%%%%%%%%%%

Define

\[
\boxed{
r
=
\frac{
\alpha\Delta t
}{
h^2
}.
}
\]

Then FTCS becomes

\[
\boxed{
u_i^{n+1}
=
r u_{i-1}^n
+
(1-2r)u_i^n
+
r u_{i+1}^n.
}
\]

This makes the stability condition easy to express.

%%%%%%%%%%%%%%%%%%%%%%%%%%%%%%%%%%%%%%%%%%%%%%%%%%%%%%%%%%%%
\subsection{FTCS Stability for Heat Equation}
\label{subsec:ftcs-heat-stability}
%%%%%%%%%%%%%%%%%%%%%%%%%%%%%%%%%%%%%%%%%%%%%%%%%%%%%%%%%%%%

For the one-dimensional heat equation, von Neumann analysis gives the
condition

\[
\boxed{
0\leq r\leq\frac12.
}
\]

Thus,

\[
\boxed{
\Delta t
\leq
\frac{
h^2
}{
2\alpha
}.
}
\]

The time step must shrink quadratically as the spatial grid is refined.

%%%%%%%%%%%%%%%%%%%%%%%%%%%%%%%%%%%%%%%%%%%%%%%%%%%%%%%%%%%%
\subsection{Worked Example: Heat Equation Stability}
\label{subsec:worked-heat-stability}
%%%%%%%%%%%%%%%%%%%%%%%%%%%%%%%%%%%%%%%%%%%%%%%%%%%%%%%%%%%%

\begin{workedexamplebox}{Choosing a Stable FTCS Time Step}

Suppose

\[
\alpha=0.5,
\]

and

\[
h=0.1.
\]

Stability requires

\[
\Delta t
\leq
\frac{
(0.1)^2
}{
2(0.5)
}.
\]

Therefore,

\[
\Delta t
\leq
0.01.
\]

\begin{answerbox}
\[
\boxed{
\Delta t\leq0.01.
}
\]
\end{answerbox}

\end{workedexamplebox}

%%%%%%%%%%%%%%%%%%%%%%%%%%%%%%%%%%%%%%%%%%%%%%%%%%%%%%%%%%%%
\subsection{Backward-Time Centered-Space Method}
\label{subsec:btcs-heat}
%%%%%%%%%%%%%%%%%%%%%%%%%%%%%%%%%%%%%%%%%%%%%%%%%%%%%%%%%%%%

Backward Euler in time gives

\[
\boxed{
\frac{
u_i^{n+1}-u_i^n
}{
\Delta t
}
=
\alpha
\frac{
u_{i+1}^{n+1}
-
2u_i^{n+1}
+
u_{i-1}^{n+1}
}{
h^2
}.
}
\]

Rearranging,

\[
\boxed{
-r u_{i-1}^{n+1}
+
(1+2r)u_i^{n+1}
-
r u_{i+1}^{n+1}
=
u_i^n.
}
\]

A linear system must be solved at every time step.

%%%%%%%%%%%%%%%%%%%%%%%%%%%%%%%%%%%%%%%%%%%%%%%%%%%%%%%%%%%%
\subsection{Backward Euler Stability for Diffusion}
\label{subsec:backward-euler-diffusion-stability}
%%%%%%%%%%%%%%%%%%%%%%%%%%%%%%%%%%%%%%%%%%%%%%%%%%%%%%%%%%%%

Backward Euler is unconditionally stable for the linear diffusion problem
in the standard von Neumann sense.

Thus stability does not impose the same restriction

\[
\Delta t
\leq
C h^2.
\]

However, a very large time step can still produce poor accuracy.

%%%%%%%%%%%%%%%%%%%%%%%%%%%%%%%%%%%%%%%%%%%%%%%%%%%%%%%%%%%%
\subsection{Crank--Nicolson Method}
\label{subsec:crank-nicolson}
%%%%%%%%%%%%%%%%%%%%%%%%%%%%%%%%%%%%%%%%%%%%%%%%%%%%%%%%%%%%

Crank--Nicolson averages the spatial operator between time levels \(n\)
and \(n+1\):

\[
\boxed{
\frac{
u_i^{n+1}-u_i^n
}{
\Delta t
}
=
\frac{\alpha}{2}
\left[
\delta_{xx}u_i^n
+
\delta_{xx}u_i^{n+1}
\right].
}
\]

Here

\[
\boxed{
\delta_{xx}u_i^n
=
\frac{
u_{i+1}^n
-
2u_i^n
+
u_{i-1}^n
}{
h^2
}.
}
\]

%%%%%%%%%%%%%%%%%%%%%%%%%%%%%%%%%%%%%%%%%%%%%%%%%%%%%%%%%%%%
\subsection{Crank--Nicolson Matrix Form}
\label{subsec:crank-nicolson-matrix}
%%%%%%%%%%%%%%%%%%%%%%%%%%%%%%%%%%%%%%%%%%%%%%%%%%%%%%%%%%%%

Let \(L_h\) denote the discrete spatial Laplacian.

Then Crank--Nicolson can be written

\[
\boxed{
\left(
I
-
\frac{
\alpha\Delta t
}{
2
}
L_h
\right)
U^{n+1}
=
\left(
I
+
\frac{
\alpha\Delta t
}{
2
}
L_h
\right)
U^n.
}
\]

The left matrix can often be factorized once and reused if the grid and
coefficients remain fixed.

%%%%%%%%%%%%%%%%%%%%%%%%%%%%%%%%%%%%%%%%%%%%%%%%%%%%%%%%%%%%
\subsection{Accuracy of Heat-Equation Schemes}
\label{subsec:heat-scheme-accuracy}
%%%%%%%%%%%%%%%%%%%%%%%%%%%%%%%%%%%%%%%%%%%%%%%%%%%%%%%%%%%%

For sufficiently smooth solutions:

\[
\boxed{
\text{FTCS}
=
O(\Delta t)
+
O(h^2),
}
\]

\[
\boxed{
\text{Backward Euler}
=
O(\Delta t)
+
O(h^2),
}
\]

and

\[
\boxed{
\text{Crank--Nicolson}
=
O(\Delta t^2)
+
O(h^2).
}
\]

Thus Crank--Nicolson is second order in both time and space.

%%%%%%%%%%%%%%%%%%%%%%%%%%%%%%%%%%%%%%%%%%%%%%%%%%%%%%%%%%%%
\subsection{Von Neumann Stability Analysis}
\label{subsec:von-neumann-analysis}
%%%%%%%%%%%%%%%%%%%%%%%%%%%%%%%%%%%%%%%%%%%%%%%%%%%%%%%%%%%%

Von Neumann analysis studies the evolution of Fourier error modes.

Assume a discrete perturbation has form

\[
\boxed{
e_i^n
=
G^n
e^{\mathrm{i}i\theta},
}
\]

where

\[
G
\]

is the amplification factor.

Substitute this mode into the numerical scheme.

Stability requires

\[
\boxed{
|G(\theta)|\leq1
}
\]

for all relevant wave numbers \(\theta\).

%%%%%%%%%%%%%%%%%%%%%%%%%%%%%%%%%%%%%%%%%%%%%%%%%%%%%%%%%%%%
\subsection{FTCS Amplification Factor}
\label{subsec:ftcs-amplification-factor}
%%%%%%%%%%%%%%%%%%%%%%%%%%%%%%%%%%%%%%%%%%%%%%%%%%%%%%%%%%%%

For the heat equation,

\[
u_i^{n+1}
=
r u_{i-1}^n
+
(1-2r)u_i^n
+
r u_{i+1}^n.
\]

Substituting

\[
e_i^n
=
G^n e^{\mathrm{i}i\theta}
\]

gives

\[
G
=
r e^{-\mathrm{i}\theta}
+
1-2r
+
r e^{\mathrm{i}\theta}.
\]

Therefore,

\[
\boxed{
G
=
1
-
4r
\sin^2
\left(
\frac{\theta}{2}
\right).
}
\]

Requiring

\[
|G|\leq1
\]

for all \(\theta\) gives

\[
\boxed{
r\leq\frac12.
}
\]

%%%%%%%%%%%%%%%%%%%%%%%%%%%%%%%%%%%%%%%%%%%%%%%%%%%%%%%%%%%%
\subsection{Advection Equation}
\label{subsec:numerical-advection-equation}
%%%%%%%%%%%%%%%%%%%%%%%%%%%%%%%%%%%%%%%%%%%%%%%%%%%%%%%%%%%%

Consider

\[
\boxed{
u_t
+
cu_x
=
0.
}
\]

For constant \(c\), the exact solution is

\[
\boxed{
u(x,t)
=
u_0(x-ct).
}
\]

Thus the initial profile is transported at speed \(c\).

A good numerical scheme should move the profile without excessive
distortion.

%%%%%%%%%%%%%%%%%%%%%%%%%%%%%%%%%%%%%%%%%%%%%%%%%%%%%%%%%%%%
\subsection{Centered FTCS for Advection}
\label{subsec:centered-ftcs-advection}
%%%%%%%%%%%%%%%%%%%%%%%%%%%%%%%%%%%%%%%%%%%%%%%%%%%%%%%%%%%%

Using forward time and centered space,

\[
\boxed{
\frac{
u_i^{n+1}-u_i^n
}{
\Delta t
}
+
c
\frac{
u_{i+1}^n-u_{i-1}^n
}{
2h
}
=
0.
}
\]

Therefore,

\[
\boxed{
u_i^{n+1}
=
u_i^n
-
\frac{\nu}{2}
\left(
u_{i+1}^n-u_{i-1}^n
\right),
}
\]

where

\[
\boxed{
\nu
=
\frac{
c\Delta t
}{
h
}.
}
\]

This standard centered FTCS scheme is unstable for pure linear advection.

%%%%%%%%%%%%%%%%%%%%%%%%%%%%%%%%%%%%%%%%%%%%%%%%%%%%%%%%%%%%
\subsection{Courant Number}
\label{subsec:courant-number}
%%%%%%%%%%%%%%%%%%%%%%%%%%%%%%%%%%%%%%%%%%%%%%%%%%%%%%%%%%%%

Define the Courant number

\[
\boxed{
\nu
=
\frac{
c\Delta t
}{
h
}.
}
\]

It measures approximately how many grid cells information travels during
one time step.

This quantity appears in stability restrictions for hyperbolic equations.

%%%%%%%%%%%%%%%%%%%%%%%%%%%%%%%%%%%%%%%%%%%%%%%%%%%%%%%%%%%%
\subsection{Upwind Method}
\label{subsec:upwind-method}
%%%%%%%%%%%%%%%%%%%%%%%%%%%%%%%%%%%%%%%%%%%%%%%%%%%%%%%%%%%%

If

\[
c>0,
\]

information travels from left to right.

Use the backward spatial difference:

\[
u_x
\approx
\frac{
u_i^n-u_{i-1}^n
}{
h
}.
\]

Then

\[
\boxed{
u_i^{n+1}
=
u_i^n
-
\nu
\left(
u_i^n-u_{i-1}^n
\right).
}
\]

Equivalently,

\[
\boxed{
u_i^{n+1}
=
(1-\nu)u_i^n
+
\nu u_{i-1}^n.
}
\]

%%%%%%%%%%%%%%%%%%%%%%%%%%%%%%%%%%%%%%%%%%%%%%%%%%%%%%%%%%%%
\subsection{Upwind Stability}
\label{subsec:upwind-stability}
%%%%%%%%%%%%%%%%%%%%%%%%%%%%%%%%%%%%%%%%%%%%%%%%%%%%%%%%%%%%

For

\[
c>0,
\]

the first-order upwind method is stable under the CFL condition

\[
\boxed{
0\leq\nu\leq1.
}
\]

Thus,

\[
\boxed{
\Delta t
\leq
\frac{h}{c}.
}
\]

For \(c<0\), the spatial difference must be reversed.

%%%%%%%%%%%%%%%%%%%%%%%%%%%%%%%%%%%%%%%%%%%%%%%%%%%%%%%%%%%%
\subsection{Direction Matters in Upwinding}
\label{subsec:upwind-direction}
%%%%%%%%%%%%%%%%%%%%%%%%%%%%%%%%%%%%%%%%%%%%%%%%%%%%%%%%%%%%

When

\[
c>0,
\]

use information from the left.

When

\[
c<0,
\]

use information from the right.

Thus upwind differentiation follows the direction from which
characteristics bring information.

%%%%%%%%%%%%%%%%%%%%%%%%%%%%%%%%%%%%%%%%%%%%%%%%%%%%%%%%%%%%
\subsection{Worked Example: Upwind Advection}
\label{subsec:worked-upwind-advection}
%%%%%%%%%%%%%%%%%%%%%%%%%%%%%%%%%%%%%%%%%%%%%%%%%%%%%%%%%%%%

\begin{workedexamplebox}{One Upwind Step}

Suppose

\[
c=1,
\qquad
h=0.2,
\qquad
\Delta t=0.1.
\]

Then

\[
\nu
=
\frac{1(0.1)}{0.2}
=
0.5.
\]

For \(c>0\),

\[
u_i^{n+1}
=
0.5u_i^n
+
0.5u_{i-1}^n.
\]

If

\[
u_i^n=4,
\]

and

\[
u_{i-1}^n=2,
\]

then

\[
u_i^{n+1}
=
0.5(4)+0.5(2)
=
3.
\]

\begin{answerbox}
\[
\boxed{
u_i^{n+1}=3.
}
\]
\end{answerbox}

\end{workedexamplebox}

%%%%%%%%%%%%%%%%%%%%%%%%%%%%%%%%%%%%%%%%%%%%%%%%%%%%%%%%%%%%
\subsection{Numerical Diffusion}
\label{subsec:numerical-diffusion}
%%%%%%%%%%%%%%%%%%%%%%%%%%%%%%%%%%%%%%%%%%%%%%%%%%%%%%%%%%%%

First-order upwind methods tend to smooth sharp solution features.

This artificial smoothing is called numerical diffusion or numerical
viscosity.

A transported discontinuity may become increasingly spread over several
grid cells.

%%%%%%%%%%%%%%%%%%%%%%%%%%%%%%%%%%%%%%%%%%%%%%%%%%%%%%%%%%%%
\subsection{Modified Equation Idea}
\label{subsec:modified-equation}
%%%%%%%%%%%%%%%%%%%%%%%%%%%%%%%%%%%%%%%%%%%%%%%%%%%%%%%%%%%%

A numerical method can often be interpreted as solving a nearby PDE.

For first-order upwind advection, the modified equation contains an
artificial diffusion term of the form

\[
\boxed{
u_t
+
cu_x
=
\varepsilon_{\mathrm{num}}u_{xx}
+
\cdots.
}
\]

The coefficient depends on the grid and time step.

Modified-equation analysis explains numerical diffusion and dispersion.

%%%%%%%%%%%%%%%%%%%%%%%%%%%%%%%%%%%%%%%%%%%%%%%%%%%%%%%%%%%%
\subsection{Numerical Dispersion}
\label{subsec:numerical-dispersion}
%%%%%%%%%%%%%%%%%%%%%%%%%%%%%%%%%%%%%%%%%%%%%%%%%%%%%%%%%%%%

Numerical dispersion occurs when different numerical wave components travel
at incorrect phase speeds.

The result may include artificial oscillations or phase errors.

This is especially important in wave propagation and high-order
hyperbolic schemes.

%%%%%%%%%%%%%%%%%%%%%%%%%%%%%%%%%%%%%%%%%%%%%%%%%%%%%%%%%%%%
\subsection{Lax--Friedrichs Method}
\label{subsec:lax-friedrichs}
%%%%%%%%%%%%%%%%%%%%%%%%%%%%%%%%%%%%%%%%%%%%%%%%%%%%%%%%%%%%

A classical scheme for linear advection is

\[
\boxed{
u_i^{n+1}
=
\frac12
\left(
u_{i+1}^n
+
u_{i-1}^n
\right)
-
\frac{\nu}{2}
\left(
u_{i+1}^n
-
u_{i-1}^n
\right).
}
\]

The averaging term adds numerical diffusion and improves stability.

%%%%%%%%%%%%%%%%%%%%%%%%%%%%%%%%%%%%%%%%%%%%%%%%%%%%%%%%%%%%
\subsection{Lax--Wendroff Method}
\label{subsec:lax-wendroff}
%%%%%%%%%%%%%%%%%%%%%%%%%%%%%%%%%%%%%%%%%%%%%%%%%%%%%%%%%%%%

For linear advection, Lax--Wendroff gives

\[
\boxed{
u_i^{n+1}
=
u_i^n
-
\frac{\nu}{2}
\left(
u_{i+1}^n-u_{i-1}^n
\right)
+
\frac{\nu^2}{2}
\left(
u_{i+1}^n
-
2u_i^n
+
u_{i-1}^n
\right).
}
\]

It is second-order accurate for smooth solutions.

Near discontinuities, however, oscillations may appear.

%%%%%%%%%%%%%%%%%%%%%%%%%%%%%%%%%%%%%%%%%%%%%%%%%%%%%%%%%%%%
\subsection{CFL Condition}
\label{subsec:cfl-condition}
%%%%%%%%%%%%%%%%%%%%%%%%%%%%%%%%%%%%%%%%%%%%%%%%%%%%%%%%%%%%

For hyperbolic equations, numerical information should propagate at least
as quickly as physical information.

This motivates the Courant--Friedrichs--Lewy condition.

Informally,

\[
\boxed{
\text{physical domain of dependence}
\subseteq
\text{numerical domain of dependence}.
}
\]

For simple explicit advection schemes, this commonly produces

\[
\boxed{
|\nu|
\leq1.
}
\]

%%%%%%%%%%%%%%%%%%%%%%%%%%%%%%%%%%%%%%%%%%%%%%%%%%%%%%%%%%%%
\subsection{Wave Equation}
\label{subsec:numerical-wave-equation}
%%%%%%%%%%%%%%%%%%%%%%%%%%%%%%%%%%%%%%%%%%%%%%%%%%%%%%%%%%%%

Consider

\[
\boxed{
u_{tt}
=
c^2u_{xx}.
}
\]

Use centered differences in both time and space:

\[
\frac{
u_i^{n+1}
-
2u_i^n
+
u_i^{n-1}
}{
\Delta t^2
}
=
c^2
\frac{
u_{i+1}^n
-
2u_i^n
+
u_{i-1}^n
}{
h^2
}.
\]

%%%%%%%%%%%%%%%%%%%%%%%%%%%%%%%%%%%%%%%%%%%%%%%%%%%%%%%%%%%%
\subsection{Explicit Wave Update}
\label{subsec:explicit-wave-update}
%%%%%%%%%%%%%%%%%%%%%%%%%%%%%%%%%%%%%%%%%%%%%%%%%%%%%%%%%%%%

Define

\[
\boxed{
\nu
=
\frac{
c\Delta t
}{
h
}.
}
\]

Then

\[
\boxed{
u_i^{n+1}
=
2u_i^n
-
u_i^{n-1}
+
\nu^2
\left(
u_{i+1}^n
-
2u_i^n
+
u_{i-1}^n
\right).
}
\]

The scheme is second-order accurate in both space and time.

%%%%%%%%%%%%%%%%%%%%%%%%%%%%%%%%%%%%%%%%%%%%%%%%%%%%%%%%%%%%
\subsection{Wave Stability Condition}
\label{subsec:wave-stability-condition}
%%%%%%%%%%%%%%%%%%%%%%%%%%%%%%%%%%%%%%%%%%%%%%%%%%%%%%%%%%%%

For the standard one-dimensional explicit centered scheme,

\[
\boxed{
|\nu|
\leq1
}
\]

is required for stability.

Thus,

\[
\boxed{
\Delta t
\leq
\frac{h}{|c|}.
}
\]

This is another CFL-type restriction.

%%%%%%%%%%%%%%%%%%%%%%%%%%%%%%%%%%%%%%%%%%%%%%%%%%%%%%%%%%%%
\subsection{Starting the Wave Scheme}
\label{subsec:starting-wave-scheme}
%%%%%%%%%%%%%%%%%%%%%%%%%%%%%%%%%%%%%%%%%%%%%%%%%%%%%%%%%%%%

The second-order time recurrence requires both

\[
u_i^0
\]

and

\[
u_i^1.
\]

If initial displacement is

\[
u(x,0)=f(x)
\]

and initial velocity is

\[
u_t(x,0)=g(x),
\]

Taylor expansion can generate \(u_i^1\).

A second-order start is

\[
\boxed{
u_i^1
=
u_i^0
+
\Delta t\,g_i
+
\frac{\nu^2}{2}
\left(
u_{i+1}^0
-
2u_i^0
+
u_{i-1}^0
\right).
}
\]

%%%%%%%%%%%%%%%%%%%%%%%%%%%%%%%%%%%%%%%%%%%%%%%%%%%%%%%%%%%%
\subsection{Conservation Laws}
\label{subsec:numerical-conservation-laws}
%%%%%%%%%%%%%%%%%%%%%%%%%%%%%%%%%%%%%%%%%%%%%%%%%%%%%%%%%%%%

A one-dimensional conservation law has form

\[
\boxed{
u_t
+
f(u)_x
=
0.
}
\]

Integrating over cell

\[
[x_{i-1/2},x_{i+1/2}]
\]

gives

\[
\boxed{
\frac{d}{dt}
\int_{x_{i-1/2}}^{x_{i+1/2}}
u(x,t)\,dx
=
-
f(u(x_{i+1/2},t))
+
f(u(x_{i-1/2},t)).
}
\]

This integral structure motivates finite-volume methods.

%%%%%%%%%%%%%%%%%%%%%%%%%%%%%%%%%%%%%%%%%%%%%%%%%%%%%%%%%%%%
\subsection{Finite-Volume Method}
\label{subsec:finite-volume-method}
%%%%%%%%%%%%%%%%%%%%%%%%%%%%%%%%%%%%%%%%%%%%%%%%%%%%%%%%%%%%

Define the cell average

\[
\boxed{
\overline{u}_i(t)
=
\frac1{\Delta x}
\int_{x_{i-1/2}}^{x_{i+1/2}}
u(x,t)\,dx.
}
\]

Then

\[
\boxed{
\frac{d\overline{u}_i}{dt}
=
-
\frac1{\Delta x}
\left[
F_{i+1/2}
-
F_{i-1/2}
\right],
}
\]

where \(F_{i\pm1/2}\) are numerical fluxes.

%%%%%%%%%%%%%%%%%%%%%%%%%%%%%%%%%%%%%%%%%%%%%%%%%%%%%%%%%%%%
\subsection{Conservation in Finite Volumes}
\label{subsec:finite-volume-conservation}
%%%%%%%%%%%%%%%%%%%%%%%%%%%%%%%%%%%%%%%%%%%%%%%%%%%%%%%%%%%%

The flux leaving one cell enters the neighboring cell.

Thus internal fluxes cancel when equations are summed over cells.

This gives a discrete conservation property.

Finite-volume methods are therefore especially natural for conservation
laws and fluid dynamics.

%%%%%%%%%%%%%%%%%%%%%%%%%%%%%%%%%%%%%%%%%%%%%%%%%%%%%%%%%%%%
\subsection{Numerical Flux}
\label{subsec:numerical-flux}
%%%%%%%%%%%%%%%%%%%%%%%%%%%%%%%%%%%%%%%%%%%%%%%%%%%%%%%%%%%%

A numerical flux approximates the physical flux at an interface:

\[
\boxed{
F_{i+1/2}
=
\widehat{f}
(
u_i,
u_{i+1}
).
}
\]

A consistent flux satisfies

\[
\boxed{
\widehat{f}(u,u)=f(u).
}
\]

Different numerical fluxes produce different stability and dissipation
properties.

%%%%%%%%%%%%%%%%%%%%%%%%%%%%%%%%%%%%%%%%%%%%%%%%%%%%%%%%%%%%
\subsection{Godunov Method Preview}
\label{subsec:godunov-method}
%%%%%%%%%%%%%%%%%%%%%%%%%%%%%%%%%%%%%%%%%%%%%%%%%%%%%%%%%%%%

Godunov's method treats neighboring cell averages as piecewise constant
states.

At each cell interface, solve or approximate a local Riemann problem.

The resulting interface solution determines the numerical flux.

This framework forms the basis of many modern high-resolution methods.

%%%%%%%%%%%%%%%%%%%%%%%%%%%%%%%%%%%%%%%%%%%%%%%%%%%%%%%%%%%%
\subsection{Shock Waves}
\label{subsec:numerical-shock-waves}
%%%%%%%%%%%%%%%%%%%%%%%%%%%%%%%%%%%%%%%%%%%%%%%%%%%%%%%%%%%%

Nonlinear hyperbolic conservation laws can develop discontinuities even
from smooth initial data.

Classical derivatives then fail at the shock.

The appropriate solution concept becomes a weak solution.

Numerical methods must also select the physically relevant entropy
solution.

%%%%%%%%%%%%%%%%%%%%%%%%%%%%%%%%%%%%%%%%%%%%%%%%%%%%%%%%%%%%
\subsection{Monotone Schemes}
\label{subsec:monotone-schemes}
%%%%%%%%%%%%%%%%%%%%%%%%%%%%%%%%%%%%%%%%%%%%%%%%%%%%%%%%%%%%

Monotone schemes prevent certain kinds of artificial oscillation by making
the update monotonic in its input values.

For scalar conservation laws, monotonicity provides strong stability and
convergence properties.

However, monotone methods are generally limited to first-order accuracy
near smooth regions under standard linearity assumptions.

%%%%%%%%%%%%%%%%%%%%%%%%%%%%%%%%%%%%%%%%%%%%%%%%%%%%%%%%%%%%
\subsection{Total Variation}
\label{subsec:total-variation}
%%%%%%%%%%%%%%%%%%%%%%%%%%%%%%%%%%%%%%%%%%%%%%%%%%%%%%%%%%%%

For grid values \(u_i\), define discrete total variation by

\[
\boxed{
TV(u)
=
\sum_i
|u_{i+1}-u_i|.
}
\]

A total-variation-diminishing method satisfies

\[
\boxed{
TV(u^{n+1})
\leq
TV(u^n).
}
\]

This helps suppress spurious oscillations near discontinuities.

%%%%%%%%%%%%%%%%%%%%%%%%%%%%%%%%%%%%%%%%%%%%%%%%%%%%%%%%%%%%
\subsection{Flux Limiters Preview}
\label{subsec:flux-limiters}
%%%%%%%%%%%%%%%%%%%%%%%%%%%%%%%%%%%%%%%%%%%%%%%%%%%%%%%%%%%%

High-resolution schemes often blend:

\[
\boxed{
\text{stable low-order fluxes}
}
\]

with

\[
\boxed{
\text{higher-order corrections}.
}
\]

Flux limiters reduce the high-order correction near discontinuities while
retaining higher accuracy in smooth regions.

%%%%%%%%%%%%%%%%%%%%%%%%%%%%%%%%%%%%%%%%%%%%%%%%%%%%%%%%%%%%
\subsection{Finite-Element Method}
\label{subsec:finite-element-method}
%%%%%%%%%%%%%%%%%%%%%%%%%%%%%%%%%%%%%%%%%%%%%%%%%%%%%%%%%%%%

Finite-element methods begin from a weak or variational formulation rather
than directly approximating derivatives pointwise.

Consider Poisson's equation

\[
\boxed{
-\Delta u=f
}
\]

on domain \(\Omega\), with

\[
u=0
\]

on the boundary.

Multiply by test function \(v\) and integrate:

\[
\boxed{
\int_\Omega
(-\Delta u)v\,dx
=
\int_\Omega
fv\,dx.
}
\]

%%%%%%%%%%%%%%%%%%%%%%%%%%%%%%%%%%%%%%%%%%%%%%%%%%%%%%%%%%%%
\subsection{Weak Form of Poisson's Equation}
\label{subsec:weak-form-poisson}
%%%%%%%%%%%%%%%%%%%%%%%%%%%%%%%%%%%%%%%%%%%%%%%%%%%%%%%%%%%%

Using integration by parts,

\[
\int_\Omega
\nabla u\cdot\nabla v\,dx
-
\int_{\partial\Omega}
\frac{\partial u}{\partial n}v\,ds
=
\int_\Omega
fv\,dx.
\]

For homogeneous Dirichlet test functions, the boundary term vanishes in
the standard weak formulation.

Thus seek \(u\) such that

\[
\boxed{
\int_\Omega
\nabla u\cdot\nabla v\,dx
=
\int_\Omega
fv\,dx
}
\]

for all admissible test functions \(v\).

%%%%%%%%%%%%%%%%%%%%%%%%%%%%%%%%%%%%%%%%%%%%%%%%%%%%%%%%%%%%
\subsection{Why Weak Formulations Matter}
\label{subsec:why-weak-form}
%%%%%%%%%%%%%%%%%%%%%%%%%%%%%%%%%%%%%%%%%%%%%%%%%%%%%%%%%%%%

The strong Poisson equation requires second derivatives of \(u\).

The weak formulation requires only first derivatives.

Thus weak formulations permit less regular solution spaces and provide a
natural foundation for finite elements.

%%%%%%%%%%%%%%%%%%%%%%%%%%%%%%%%%%%%%%%%%%%%%%%%%%%%%%%%%%%%
\subsection{Trial and Test Spaces}
\label{subsec:trial-test-spaces}
%%%%%%%%%%%%%%%%%%%%%%%%%%%%%%%%%%%%%%%%%%%%%%%%%%%%%%%%%%%%

Let

\[
V
\]

be an appropriate function space.

The exact weak problem is:

find

\[
\boxed{
u\in V
}
\]

such that

\[
\boxed{
a(u,v)
=
\ell(v)
}
\]

for every

\[
\boxed{
v\in V.
}
\]

Here,

\[
a(u,v)
\]

is a bilinear form and

\[
\ell(v)
\]

is a linear functional.

%%%%%%%%%%%%%%%%%%%%%%%%%%%%%%%%%%%%%%%%%%%%%%%%%%%%%%%%%%%%
\subsection{Finite-Dimensional Approximation Space}
\label{subsec:finite-element-space}
%%%%%%%%%%%%%%%%%%%%%%%%%%%%%%%%%%%%%%%%%%%%%%%%%%%%%%%%%%%%

Choose a finite-dimensional subspace

\[
\boxed{
V_h
\subset V.
}
\]

Let basis functions be

\[
\boxed{
\phi_1,\ldots,\phi_N.
}
\]

Approximate

\[
\boxed{
u_h
=
\sum_{j=1}^{N}
U_j\phi_j.
}
\]

The unknowns are now the coefficients

\[
U_j.
\]

%%%%%%%%%%%%%%%%%%%%%%%%%%%%%%%%%%%%%%%%%%%%%%%%%%%%%%%%%%%%
\subsection{Galerkin Method}
\label{subsec:galerkin-method}
%%%%%%%%%%%%%%%%%%%%%%%%%%%%%%%%%%%%%%%%%%%%%%%%%%%%%%%%%%%%

The Galerkin method uses the same finite-dimensional space for trial and
test functions.

Require

\[
\boxed{
a(u_h,v_h)
=
\ell(v_h)
}
\]

for every

\[
v_h\in V_h.
\]

Choosing

\[
v_h=\phi_i
\]

gives

\[
\boxed{
\sum_{j=1}^{N}
U_j
a(\phi_j,\phi_i)
=
\ell(\phi_i).
}
\]

%%%%%%%%%%%%%%%%%%%%%%%%%%%%%%%%%%%%%%%%%%%%%%%%%%%%%%%%%%%%
\subsection{Finite-Element Matrix System}
\label{subsec:finite-element-matrix}
%%%%%%%%%%%%%%%%%%%%%%%%%%%%%%%%%%%%%%%%%%%%%%%%%%%%%%%%%%%%

Define

\[
\boxed{
K_{ij}
=
a(\phi_j,\phi_i),
}
\]

and

\[
\boxed{
F_i
=
\ell(\phi_i).
}
\]

Then

\[
\boxed{
KU=F.
}
\]

For Poisson's equation,

\[
\boxed{
K_{ij}
=
\int_\Omega
\nabla\phi_j
\cdot
\nabla\phi_i
\,dx.
}
\]

The matrix \(K\) is called the stiffness matrix.

%%%%%%%%%%%%%%%%%%%%%%%%%%%%%%%%%%%%%%%%%%%%%%%%%%%%%%%%%%%%
\subsection{Mass Matrix}
\label{subsec:finite-element-mass-matrix}
%%%%%%%%%%%%%%%%%%%%%%%%%%%%%%%%%%%%%%%%%%%%%%%%%%%%%%%%%%%%

For time-dependent problems, terms involving

\[
u_t
\]

often produce the mass matrix

\[
\boxed{
M_{ij}
=
\int_\Omega
\phi_j\phi_i
\,dx.
}
\]

After spatial discretization, a PDE may become

\[
\boxed{
M\dot{U}
+
KU
=
F.
}
\]

This is a system of ODEs or differential-algebraic equations.

%%%%%%%%%%%%%%%%%%%%%%%%%%%%%%%%%%%%%%%%%%%%%%%%%%%%%%%%%%%%
\subsection{Piecewise Linear Finite Elements}
\label{subsec:piecewise-linear-fem}
%%%%%%%%%%%%%%%%%%%%%%%%%%%%%%%%%%%%%%%%%%%%%%%%%%%%%%%%%%%%

In one dimension, a common basis uses piecewise linear hat functions.

Each basis function satisfies

\[
\boxed{
\phi_i(x_j)
=
\delta_{ij}.
}
\]

It is nonzero only on intervals adjacent to node \(x_i\).

This local support produces sparse finite-element matrices.

%%%%%%%%%%%%%%%%%%%%%%%%%%%%%%%%%%%%%%%%%%%%%%%%%%%%%%%%%%%%
\subsection{Local Element Matrices}
\label{subsec:local-element-matrices}
%%%%%%%%%%%%%%%%%%%%%%%%%%%%%%%%%%%%%%%%%%%%%%%%%%%%%%%%%%%%

Finite-element calculations are usually performed element by element.

For each element \(e\), compute a local stiffness matrix

\[
\boxed{
K^{(e)}.
}
\]

Then insert its entries into the appropriate locations of the global
matrix.

This process is called assembly.

%%%%%%%%%%%%%%%%%%%%%%%%%%%%%%%%%%%%%%%%%%%%%%%%%%%%%%%%%%%%
\subsection{Finite-Element Assembly}
\label{subsec:finite-element-assembly}
%%%%%%%%%%%%%%%%%%%%%%%%%%%%%%%%%%%%%%%%%%%%%%%%%%%%%%%%%%%%

Because neighboring elements share nodes, their local contributions add.

Symbolically,

\[
\boxed{
K
=
\sum_e
A_e^T
K^{(e)}
A_e,
}
\]

where \(A_e\) represents the mapping between local and global degrees of
freedom.

Assembly is a central finite-element operation.

%%%%%%%%%%%%%%%%%%%%%%%%%%%%%%%%%%%%%%%%%%%%%%%%%%%%%%%%%%%%
\subsection{Advantages of Finite Elements}
\label{subsec:fem-advantages}
%%%%%%%%%%%%%%%%%%%%%%%%%%%%%%%%%%%%%%%%%%%%%%%%%%%%%%%%%%%%

Finite elements are especially useful for:

\[
\boxed{
\text{irregular geometries},
}
\]

\[
\boxed{
\text{variable coefficients},
}
\]

\[
\boxed{
\text{unstructured meshes},
}
\]

and

\[
\boxed{
\text{complex boundary conditions}.
}
\]

They also have a strong variational mathematical foundation.

%%%%%%%%%%%%%%%%%%%%%%%%%%%%%%%%%%%%%%%%%%%%%%%%%%%%%%%%%%%%
\subsection{Structured and Unstructured Meshes}
\label{subsec:structured-unstructured-meshes}
%%%%%%%%%%%%%%%%%%%%%%%%%%%%%%%%%%%%%%%%%%%%%%%%%%%%%%%%%%%%

A structured mesh has a regular indexing pattern.

Examples include Cartesian grids.

An unstructured mesh consists of irregularly connected elements such as:

\[
\boxed{
\text{triangles},
}
\]

\[
\boxed{
\text{tetrahedra},
}
\]

or other polygons and polyhedra.

Unstructured meshes fit complex geometry more naturally.

%%%%%%%%%%%%%%%%%%%%%%%%%%%%%%%%%%%%%%%%%%%%%%%%%%%%%%%%%%%%
\subsection{Mesh Size}
\label{subsec:pde-mesh-size}
%%%%%%%%%%%%%%%%%%%%%%%%%%%%%%%%%%%%%%%%%%%%%%%%%%%%%%%%%%%%

The parameter

\[
\boxed{
h
}
\]

often denotes a characteristic element diameter or grid spacing.

Smaller \(h\) generally improves spatial resolution.

But reducing \(h\) increases:

\[
\boxed{
\text{degrees of freedom},
}
\]

\[
\boxed{
\text{memory},
}
\]

and

\[
\boxed{
\text{computational cost}.
}
\]

%%%%%%%%%%%%%%%%%%%%%%%%%%%%%%%%%%%%%%%%%%%%%%%%%%%%%%%%%%%%
\subsection{Adaptive Mesh Refinement}
\label{subsec:adaptive-mesh-refinement}
%%%%%%%%%%%%%%%%%%%%%%%%%%%%%%%%%%%%%%%%%%%%%%%%%%%%%%%%%%%%

Uniform refinement adds resolution everywhere.

Adaptive mesh refinement, abbreviated AMR, places more resolution only
where needed.

A typical procedure is:

\begin{enumerate}
    \item solve on the current mesh;
    \item estimate local error;
    \item mark high-error regions;
    \item refine those regions;
    \item solve again.
\end{enumerate}

%%%%%%%%%%%%%%%%%%%%%%%%%%%%%%%%%%%%%%%%%%%%%%%%%%%%%%%%%%%%
\subsection{\(h\)- and \(p\)-Refinement}
\label{subsec:h-p-refinement}
%%%%%%%%%%%%%%%%%%%%%%%%%%%%%%%%%%%%%%%%%%%%%%%%%%%%%%%%%%%%

In finite-element terminology:

\[
\boxed{
h\text{-refinement}
=
\text{smaller elements},
}
\]

while

\[
\boxed{
p\text{-refinement}
=
\text{higher polynomial degree}.
}
\]

Combining both gives

\[
\boxed{
hp\text{-methods}.
}
\]

These can achieve very high accuracy for sufficiently smooth solutions.

%%%%%%%%%%%%%%%%%%%%%%%%%%%%%%%%%%%%%%%%%%%%%%%%%%%%%%%%%%%%
\subsection{Spectral Methods}
\label{subsec:spectral-methods}
%%%%%%%%%%%%%%%%%%%%%%%%%%%%%%%%%%%%%%%%%%%%%%%%%%%%%%%%%%%%

Spectral methods approximate a solution using global basis functions:

\[
\boxed{
u_N(x)
=
\sum_{k=0}^{N}
a_k\phi_k(x).
}
\]

Common bases include:

\[
\boxed{
\text{Fourier functions},
}
\]

\[
\boxed{
\text{Chebyshev polynomials},
}
\]

and

\[
\boxed{
\text{Legendre polynomials}.
}
\]

%%%%%%%%%%%%%%%%%%%%%%%%%%%%%%%%%%%%%%%%%%%%%%%%%%%%%%%%%%%%
\subsection{Spectral Accuracy}
\label{subsec:spectral-accuracy}
%%%%%%%%%%%%%%%%%%%%%%%%%%%%%%%%%%%%%%%%%%%%%%%%%%%%%%%%%%%%

For very smooth or analytic solutions, spectral methods can converge
extremely rapidly.

The error may decrease faster than any fixed algebraic power of

\[
1/N.
\]

This is sometimes called spectral or exponential convergence, depending on
the precise regularity and setting.

%%%%%%%%%%%%%%%%%%%%%%%%%%%%%%%%%%%%%%%%%%%%%%%%%%%%%%%%%%%%
\subsection{Fourier Spectral Method}
\label{subsec:fourier-spectral-method}
%%%%%%%%%%%%%%%%%%%%%%%%%%%%%%%%%%%%%%%%%%%%%%%%%%%%%%%%%%%%

For periodic problems, approximate

\[
u(x)
\]

by a Fourier series:

\[
\boxed{
u_N(x)
=
\sum_{k=-N/2}^{N/2}
\widehat{u}_k
e^{\mathrm{i}kx}.
}
\]

Differentiation becomes multiplication in Fourier space:

\[
\boxed{
\widehat{u_x}_k
=
\mathrm{i}k\widehat{u}_k.
}
\]

This can make high-order spatial derivatives highly accurate.

%%%%%%%%%%%%%%%%%%%%%%%%%%%%%%%%%%%%%%%%%%%%%%%%%%%%%%%%%%%%
\subsection{Pseudospectral Methods}
\label{subsec:pseudospectral-methods}
%%%%%%%%%%%%%%%%%%%%%%%%%%%%%%%%%%%%%%%%%%%%%%%%%%%%%%%%%%%%

Pseudospectral methods represent the solution globally but evaluate
nonlinear terms at selected collocation points.

Derivatives may be computed using differentiation matrices:

\[
\boxed{
U_x
\approx
DU.
}
\]

This combines spectral approximation with nodal computation.

%%%%%%%%%%%%%%%%%%%%%%%%%%%%%%%%%%%%%%%%%%%%%%%%%%%%%%%%%%%%
\subsection{Spectral Limitations}
\label{subsec:spectral-limitations}
%%%%%%%%%%%%%%%%%%%%%%%%%%%%%%%%%%%%%%%%%%%%%%%%%%%%%%%%%%%%

Global spectral methods can be less convenient for:

\[
\boxed{
\text{complex geometry},
}
\]

\[
\boxed{
\text{discontinuous solutions},
}
\]

or

\[
\boxed{
\text{localized nonsmooth features}.
}
\]

Discontinuities can produce Gibbs oscillations.

%%%%%%%%%%%%%%%%%%%%%%%%%%%%%%%%%%%%%%%%%%%%%%%%%%%%%%%%%%%%
\subsection{Method of Lines}
\label{subsec:pde-method-of-lines}
%%%%%%%%%%%%%%%%%%%%%%%%%%%%%%%%%%%%%%%%%%%%%%%%%%%%%%%%%%%%

The method of lines discretizes spatial derivatives but leaves time
continuous.

For a PDE

\[
u_t
=
\mathcal{L}u,
\]

spatial discretization gives

\[
\boxed{
\frac{dU}{dt}
=
L_hU.
}
\]

This large ODE system can then be integrated using the methods of
Section~\ref{sec:numerical-odes}.

%%%%%%%%%%%%%%%%%%%%%%%%%%%%%%%%%%%%%%%%%%%%%%%%%%%%%%%%%%%%
\subsection{Semi-Discrete Heat Equation}
\label{subsec:semidiscrete-heat}
%%%%%%%%%%%%%%%%%%%%%%%%%%%%%%%%%%%%%%%%%%%%%%%%%%%%%%%%%%%%

For

\[
u_t
=
\alpha u_{xx},
\]

centered spatial differences give

\[
\boxed{
\frac{dU}{dt}
=
\alpha L_hU,
}
\]

where \(L_h\) is the discrete Laplacian.

The resulting ODE system can be stiff because eigenvalues of \(L_h\) grow
in magnitude as

\[
h\to0.
\]

%%%%%%%%%%%%%%%%%%%%%%%%%%%%%%%%%%%%%%%%%%%%%%%%%%%%%%%%%%%%
\subsection{Why PDE Semi-Discretizations Become Stiff}
\label{subsec:pde-semidiscrete-stiffness}
%%%%%%%%%%%%%%%%%%%%%%%%%%%%%%%%%%%%%%%%%%%%%%%%%%%%%%%%%%%%

For diffusion, the largest-magnitude eigenvalues of the discrete Laplacian
scale approximately like

\[
\boxed{
O(h^{-2}).
}
\]

As the mesh is refined, fast-decaying numerical modes become increasingly
rapid.

This explains the explicit restriction

\[
\boxed{
\Delta t=O(h^2)
}
\]

for diffusion problems.

%%%%%%%%%%%%%%%%%%%%%%%%%%%%%%%%%%%%%%%%%%%%%%%%%%%%%%%%%%%%
\subsection{Operator Splitting}
\label{subsec:operator-splitting}
%%%%%%%%%%%%%%%%%%%%%%%%%%%%%%%%%%%%%%%%%%%%%%%%%%%%%%%%%%%%

Suppose

\[
\boxed{
u_t
=
(A+B)u.
}
\]

Instead of solving the full operator simultaneously, splitting methods
alternate between subproblems involving \(A\) and \(B\).

For one step, a first-order splitting is conceptually:

\[
\boxed{
u^*
=
\text{evolve under }A,
}
\]

then

\[
\boxed{
u^{n+1}
=
\text{evolve under }B.
}
\]

%%%%%%%%%%%%%%%%%%%%%%%%%%%%%%%%%%%%%%%%%%%%%%%%%%%%%%%%%%%%
\subsection{Strang Splitting}
\label{subsec:strang-splitting}
%%%%%%%%%%%%%%%%%%%%%%%%%%%%%%%%%%%%%%%%%%%%%%%%%%%%%%%%%%%%

A symmetric second-order splitting is:

\[
\boxed{
A
\text{ for half a step},
}
\]

then

\[
\boxed{
B
\text{ for a full step},
}
\]

then

\[
\boxed{
A
\text{ for half a step}.
}
\]

This is known as Strang splitting.

It is widely used in reaction-diffusion and multiphysics problems.

%%%%%%%%%%%%%%%%%%%%%%%%%%%%%%%%%%%%%%%%%%%%%%%%%%%%%%%%%%%%
\subsection{Nonlinear PDEs}
\label{subsec:numerical-nonlinear-pdes}
%%%%%%%%%%%%%%%%%%%%%%%%%%%%%%%%%%%%%%%%%%%%%%%%%%%%%%%%%%%%

A nonlinear PDE may produce a nonlinear algebraic system after
discretization:

\[
\boxed{
F(U)=0.
}
\]

Examples include:

\[
\boxed{
\text{nonlinear diffusion},
}
\]

\[
\boxed{
\text{Navier--Stokes equations},
}
\]

and

\[
\boxed{
\text{nonlinear reaction-diffusion systems}.
}
\]

%%%%%%%%%%%%%%%%%%%%%%%%%%%%%%%%%%%%%%%%%%%%%%%%%%%%%%%%%%%%
\subsection{Newton Linearization for PDE Systems}
\label{subsec:newton-linearization-pde}
%%%%%%%%%%%%%%%%%%%%%%%%%%%%%%%%%%%%%%%%%%%%%%%%%%%%%%%%%%%%

Given

\[
F(U)=0,
\]

Newton's method solves

\[
\boxed{
J_F(U^{(k)})
\Delta U
=
-
F(U^{(k)}),
}
\]

then updates

\[
\boxed{
U^{(k+1)}
=
U^{(k)}
+
\Delta U.
}
\]

Each Newton iteration therefore requires a large linear solve.

%%%%%%%%%%%%%%%%%%%%%%%%%%%%%%%%%%%%%%%%%%%%%%%%%%%%%%%%%%%%
\subsection{Picard Iteration}
\label{subsec:picard-pde}
%%%%%%%%%%%%%%%%%%%%%%%%%%%%%%%%%%%%%%%%%%%%%%%%%%%%%%%%%%%%

A simpler nonlinear iteration may lag nonlinear coefficients.

For example, if

\[
-\nabla\cdot
(
a(u)\nabla u
)
=
f,
\]

one may solve

\[
\boxed{
-\nabla\cdot
(
a(u^{(k)})\nabla u^{(k+1)}
)
=
f.
}
\]

This is a fixed-point or Picard-type iteration.

It may converge more slowly than Newton but can be simpler and more
robust.

%%%%%%%%%%%%%%%%%%%%%%%%%%%%%%%%%%%%%%%%%%%%%%%%%%%%%%%%%%%%
\subsection{Linear Solvers for PDE Systems}
\label{subsec:pde-linear-solvers}
%%%%%%%%%%%%%%%%%%%%%%%%%%%%%%%%%%%%%%%%%%%%%%%%%%%%%%%%%%%%

PDE matrices can contain millions or billions of unknowns.

Direct factorization may become too expensive.

Large-scale solvers often use:

\[
\boxed{
\text{conjugate gradient},
}
\]

\[
\boxed{
\text{GMRES},
}
\]

or related Krylov methods.

Preconditioning is usually essential.

%%%%%%%%%%%%%%%%%%%%%%%%%%%%%%%%%%%%%%%%%%%%%%%%%%%%%%%%%%%%
\subsection{Multigrid Idea}
\label{subsec:multigrid}
%%%%%%%%%%%%%%%%%%%%%%%%%%%%%%%%%%%%%%%%%%%%%%%%%%%%%%%%%%%%

Iterative methods often reduce oscillatory errors quickly but struggle with
smooth long-wavelength errors.

On a coarser grid, those smooth errors appear more oscillatory and can be
removed efficiently.

Multigrid combines:

\[
\boxed{
\text{smoothing on fine grids}
}
\]

with

\[
\boxed{
\text{correction on coarse grids}.
}
\]

%%%%%%%%%%%%%%%%%%%%%%%%%%%%%%%%%%%%%%%%%%%%%%%%%%%%%%%%%%%%
\subsection{Two-Grid Correction}
\label{subsec:two-grid-correction}
%%%%%%%%%%%%%%%%%%%%%%%%%%%%%%%%%%%%%%%%%%%%%%%%%%%%%%%%%%%%

A basic two-grid cycle is:

\begin{enumerate}
    \item smooth the fine-grid approximation;
    \item compute the residual;
    \item restrict the residual to a coarser grid;
    \item solve or approximate the coarse error equation;
    \item interpolate the correction to the fine grid;
    \item correct the fine-grid solution;
    \item smooth again.
\end{enumerate}

This idea extends recursively to multigrid.

%%%%%%%%%%%%%%%%%%%%%%%%%%%%%%%%%%%%%%%%%%%%%%%%%%%%%%%%%%%%
\subsection{Why Multigrid Is Powerful}
\label{subsec:why-multigrid}
%%%%%%%%%%%%%%%%%%%%%%%%%%%%%%%%%%%%%%%%%%%%%%%%%%%%%%%%%%%%

For many elliptic problems, carefully designed multigrid methods can achieve
computational work approximately proportional to the number of unknowns:

\[
\boxed{
O(N).
}
\]

This is close to optimal because simply reading the \(N\) unknowns already
requires \(O(N)\) work.

%%%%%%%%%%%%%%%%%%%%%%%%%%%%%%%%%%%%%%%%%%%%%%%%%%%%%%%%%%%%
\subsection{Preconditioning PDE Systems}
\label{subsec:pde-preconditioning}
%%%%%%%%%%%%%%%%%%%%%%%%%%%%%%%%%%%%%%%%%%%%%%%%%%%%%%%%%%%%

Useful preconditioners include:

\[
\boxed{
\text{Jacobi},
}
\]

\[
\boxed{
\text{incomplete factorization},
}
\]

\[
\boxed{
\text{domain decomposition},
}
\]

and

\[
\boxed{
\text{multigrid}.
}
\]

The best choice depends strongly on the PDE and discretization.

%%%%%%%%%%%%%%%%%%%%%%%%%%%%%%%%%%%%%%%%%%%%%%%%%%%%%%%%%%%%
\subsection{Domain Decomposition}
\label{subsec:domain-decomposition}
%%%%%%%%%%%%%%%%%%%%%%%%%%%%%%%%%%%%%%%%%%%%%%%%%%%%%%%%%%%%

A large physical domain can be divided into subdomains.

Each processor or solver handles a local problem.

Information is exchanged across subdomain boundaries.

This provides both:

\[
\boxed{
\text{preconditioning}
}
\]

and

\[
\boxed{
\text{parallelism}.
}
\]

%%%%%%%%%%%%%%%%%%%%%%%%%%%%%%%%%%%%%%%%%%%%%%%%%%%%%%%%%%%%
\subsection{Consistency of PDE Schemes}
\label{subsec:pde-consistency}
%%%%%%%%%%%%%%%%%%%%%%%%%%%%%%%%%%%%%%%%%%%%%%%%%%%%%%%%%%%%

A finite-difference scheme is consistent if its discrete operator
approaches the differential operator as

\[
h,\Delta t\to0.
\]

For example,

\[
\boxed{
\frac{
u_{i+1}-2u_i+u_{i-1}
}{
h^2
}
=
u_{xx}(x_i)
+
O(h^2).
}
\]

Thus the centered second derivative is second-order consistent.

%%%%%%%%%%%%%%%%%%%%%%%%%%%%%%%%%%%%%%%%%%%%%%%%%%%%%%%%%%%%
\subsection{PDE Stability}
\label{subsec:pde-stability}
%%%%%%%%%%%%%%%%%%%%%%%%%%%%%%%%%%%%%%%%%%%%%%%%%%%%%%%%%%%%

PDE stability describes whether perturbations remain controlled under the
discrete evolution or algebraic solve.

The appropriate stability notion depends on the PDE class and
discretization.

Common tools include:

\[
\boxed{
\text{von Neumann analysis},
}
\]

\[
\boxed{
\text{energy estimates},
}
\]

and

\[
\boxed{
\text{matrix norm estimates}.
}
\]

%%%%%%%%%%%%%%%%%%%%%%%%%%%%%%%%%%%%%%%%%%%%%%%%%%%%%%%%%%%%
\subsection{PDE Convergence}
\label{subsec:pde-convergence}
%%%%%%%%%%%%%%%%%%%%%%%%%%%%%%%%%%%%%%%%%%%%%%%%%%%%%%%%%%%%

A numerical PDE approximation converges if

\[
\boxed{
u_h
\to
u
}
\]

in an appropriate norm as the grid is refined.

For time-dependent schemes, both

\[
h
\]

and

\[
\Delta t
\]

may need to approach zero in a compatible way.

%%%%%%%%%%%%%%%%%%%%%%%%%%%%%%%%%%%%%%%%%%%%%%%%%%%%%%%%%%%%
\subsection{Lax Equivalence Principle Preview}
\label{subsec:lax-equivalence}
%%%%%%%%%%%%%%%%%%%%%%%%%%%%%%%%%%%%%%%%%%%%%%%%%%%%%%%%%%%%

For a well-posed linear initial-value problem and a consistent finite-
difference approximation, a classical principle states:

\[
\boxed{
\text{stability}
\Longleftrightarrow
\text{convergence}.
}
\]

This is the Lax equivalence theorem in its appropriate setting.

It does not automatically apply to arbitrary nonlinear PDEs or arbitrary
discretizations.

%%%%%%%%%%%%%%%%%%%%%%%%%%%%%%%%%%%%%%%%%%%%%%%%%%%%%%%%%%%%
\subsection{Grid Refinement Study}
\label{subsec:pde-grid-refinement}
%%%%%%%%%%%%%%%%%%%%%%%%%%%%%%%%%%%%%%%%%%%%%%%%%%%%%%%%%%%%

To verify spatial convergence:

\begin{enumerate}
    \item solve on mesh size \(h\);
    \item solve on \(h/2\);
    \item solve on \(h/4\);
    \item compare with an exact or highly resolved reference;
    \item estimate observed convergence order.
\end{enumerate}

If

\[
E(h)\approx Ch^p,
\]

then

\[
\boxed{
p
\approx
\log_2
\left(
\frac{
E(h)
}{
E(h/2)
}
\right).
}
\]

%%%%%%%%%%%%%%%%%%%%%%%%%%%%%%%%%%%%%%%%%%%%%%%%%%%%%%%%%%%%
\subsection{Space-Time Refinement}
\label{subsec:space-time-refinement}
%%%%%%%%%%%%%%%%%%%%%%%%%%%%%%%%%%%%%%%%%%%%%%%%%%%%%%%%%%%%

For a scheme with error

\[
\boxed{
O(\Delta t^q)
+
O(h^p),
}
\]

refining only one variable may eventually reveal error from the other.

A balanced convergence study should vary both:

\[
\boxed{
h
}
\]

and

\[
\boxed{
\Delta t.
}
\]

%%%%%%%%%%%%%%%%%%%%%%%%%%%%%%%%%%%%%%%%%%%%%%%%%%%%%%%%%%%%
\subsection{Manufactured Solutions}
\label{subsec:manufactured-solutions}
%%%%%%%%%%%%%%%%%%%%%%%%%%%%%%%%%%%%%%%%%%%%%%%%%%%%%%%%%%%%

When no exact benchmark solution is available, choose an artificial smooth
function

\[
u_{\mathrm{exact}}.
\]

Substitute it into the PDE to determine the forcing term and boundary
conditions that make it exact.

Then test the numerical method against this known solution.

This is called the method of manufactured solutions.

%%%%%%%%%%%%%%%%%%%%%%%%%%%%%%%%%%%%%%%%%%%%%%%%%%%%%%%%%%%%
\subsection{Verification Versus Validation}
\label{subsec:pde-verification-validation}
%%%%%%%%%%%%%%%%%%%%%%%%%%%%%%%%%%%%%%%%%%%%%%%%%%%%%%%%%%%%

Verification asks:

\[
\boxed{
\text{Are the PDE equations being solved numerically correctly?}
}
\]

Validation asks:

\[
\boxed{
\text{Does the PDE model represent the physical system adequately?}
}
\]

Grid convergence verifies the numerical method.

Experimental comparison contributes to model validation.

%%%%%%%%%%%%%%%%%%%%%%%%%%%%%%%%%%%%%%%%%%%%%%%%%%%%%%%%%%%%
\subsection{Boundary Errors}
\label{subsec:pde-boundary-errors}
%%%%%%%%%%%%%%%%%%%%%%%%%%%%%%%%%%%%%%%%%%%%%%%%%%%%%%%%%%%%

A high-order interior scheme can lose overall accuracy if low-order
boundary approximations are used.

Thus boundary discretization must be designed consistently with the desired
global accuracy.

Boundary conditions are not a minor implementation detail.

%%%%%%%%%%%%%%%%%%%%%%%%%%%%%%%%%%%%%%%%%%%%%%%%%%%%%%%%%%%%
\subsection{Ghost Points}
\label{subsec:pde-ghost-points}
%%%%%%%%%%%%%%%%%%%%%%%%%%%%%%%%%%%%%%%%%%%%%%%%%%%%%%%%%%%%

A stencil near a boundary may require points outside the physical domain.

Artificial external values are called ghost points.

Boundary conditions are used to eliminate or determine the ghost values.

Ghost-point methods can preserve centered interior stencils near
boundaries.

%%%%%%%%%%%%%%%%%%%%%%%%%%%%%%%%%%%%%%%%%%%%%%%%%%%%%%%%%%%%
\subsection{Nonuniform Grids}
\label{subsec:pde-nonuniform-grids}
%%%%%%%%%%%%%%%%%%%%%%%%%%%%%%%%%%%%%%%%%%%%%%%%%%%%%%%%%%%%

Uniform grids use constant spacing.

Nonuniform grids allow

\[
\boxed{
h_i
=
x_{i+1}-x_i
}
\]

to vary.

This permits higher resolution in selected regions.

Finite-difference coefficients must then account for unequal spacing.

%%%%%%%%%%%%%%%%%%%%%%%%%%%%%%%%%%%%%%%%%%%%%%%%%%%%%%%%%%%%
\subsection{Anisotropic Meshes}
\label{subsec:anisotropic-meshes}
%%%%%%%%%%%%%%%%%%%%%%%%%%%%%%%%%%%%%%%%%%%%%%%%%%%%%%%%%%%%

Some solutions vary rapidly in one direction but slowly in another.

An anisotropic mesh uses different scales in different directions.

This can efficiently resolve:

\[
\boxed{
\text{boundary layers},
}
\]

\[
\boxed{
\text{shocks},
}
\]

or

\[
\boxed{
\text{thin structures}.
}
\]

%%%%%%%%%%%%%%%%%%%%%%%%%%%%%%%%%%%%%%%%%%%%%%%%%%%%%%%%%%%%
\subsection{Boundary Layers}
\label{subsec:numerical-boundary-layers}
%%%%%%%%%%%%%%%%%%%%%%%%%%%%%%%%%%%%%%%%%%%%%%%%%%%%%%%%%%%%

Singularly perturbed PDEs may develop narrow regions with very large
gradients.

A uniform coarse grid can miss these layers entirely.

Possible remedies include:

\[
\boxed{
\text{mesh refinement},
}
\]

\[
\boxed{
\text{fitted meshes},
}
\]

or

\[
\boxed{
\text{stabilized methods}.
}
\]

%%%%%%%%%%%%%%%%%%%%%%%%%%%%%%%%%%%%%%%%%%%%%%%%%%%%%%%%%%%%
\subsection{Peclet Number Preview}
\label{subsec:peclet-number}
%%%%%%%%%%%%%%%%%%%%%%%%%%%%%%%%%%%%%%%%%%%%%%%%%%%%%%%%%%%%

For advection-diffusion equations, the relative strength of transport and
diffusion can be measured by a Peclet-type number:

\[
\boxed{
Pe
=
\frac{
\text{advection scale}
}{
\text{diffusion scale}
}.
}
\]

Large \(Pe\) indicates advection-dominated behavior.

Standard centered schemes may oscillate unless the grid is sufficiently
fine or stabilization is added.

%%%%%%%%%%%%%%%%%%%%%%%%%%%%%%%%%%%%%%%%%%%%%%%%%%%%%%%%%%%%
\subsection{Advection-Diffusion Equation}
\label{subsec:advection-diffusion}
%%%%%%%%%%%%%%%%%%%%%%%%%%%%%%%%%%%%%%%%%%%%%%%%%%%%%%%%%%%%

A common model is

\[
\boxed{
u_t
+
cu_x
=
\alpha u_{xx}.
}
\]

The advection term transports the solution.

The diffusion term smooths it.

A numerical method must balance the different stability and accuracy
requirements of both processes.

%%%%%%%%%%%%%%%%%%%%%%%%%%%%%%%%%%%%%%%%%%%%%%%%%%%%%%%%%%%%
\subsection{Reaction-Diffusion Equations}
\label{subsec:reaction-diffusion}
%%%%%%%%%%%%%%%%%%%%%%%%%%%%%%%%%%%%%%%%%%%%%%%%%%%%%%%%%%%%

A reaction-diffusion system has form

\[
\boxed{
u_t
=
D\Delta u
+
R(u).
}
\]

The diffusion operator couples neighboring spatial points.

The reaction term acts locally.

Operator splitting is especially natural for such systems.

%%%%%%%%%%%%%%%%%%%%%%%%%%%%%%%%%%%%%%%%%%%%%%%%%%%%%%%%%%%%
\subsection{Elliptic, Parabolic, and Hyperbolic Numerical Behavior}
\label{subsec:pde-class-behavior}
%%%%%%%%%%%%%%%%%%%%%%%%%%%%%%%%%%%%%%%%%%%%%%%%%%%%%%%%%%%%

Elliptic problems usually require solving a global boundary-value system.

Parabolic problems combine spatial coupling with one-way time evolution.

Hyperbolic problems propagate waves and information along characteristic
directions.

Thus numerical methods should reflect the mathematical structure of the
PDE class.

%%%%%%%%%%%%%%%%%%%%%%%%%%%%%%%%%%%%%%%%%%%%%%%%%%%%%%%%%%%%
\subsection{Finite Difference Versus Finite Volume}
\label{subsec:fd-vs-fv}
%%%%%%%%%%%%%%%%%%%%%%%%%%%%%%%%%%%%%%%%%%%%%%%%%%%%%%%%%%%%

Finite differences approximate derivatives directly at grid points.

Finite volumes approximate integral conservation laws over cells.

Thus:

\[
\boxed{
\text{finite difference}
\rightarrow
\text{derivative approximation},
}
\]

while

\[
\boxed{
\text{finite volume}
\rightarrow
\text{flux conservation}.
}
\]

%%%%%%%%%%%%%%%%%%%%%%%%%%%%%%%%%%%%%%%%%%%%%%%%%%%%%%%%%%%%
\subsection{Finite Difference Versus Finite Element}
\label{subsec:fd-vs-fem}
%%%%%%%%%%%%%%%%%%%%%%%%%%%%%%%%%%%%%%%%%%%%%%%%%%%%%%%%%%%%

Finite differences are especially simple on regular grids.

Finite elements are especially flexible for irregular geometry and
variational problems.

Finite differences begin from the strong differential form.

Finite elements commonly begin from a weak form.

%%%%%%%%%%%%%%%%%%%%%%%%%%%%%%%%%%%%%%%%%%%%%%%%%%%%%%%%%%%%
\subsection{Finite Element Versus Spectral}
\label{subsec:fem-vs-spectral}
%%%%%%%%%%%%%%%%%%%%%%%%%%%%%%%%%%%%%%%%%%%%%%%%%%%%%%%%%%%%

Finite elements use local basis functions.

Spectral methods often use global high-order basis functions.

Local bases are flexible for geometry and adaptivity.

Global spectral bases can provide extremely rapid convergence for smooth
solutions on suitable domains.

%%%%%%%%%%%%%%%%%%%%%%%%%%%%%%%%%%%%%%%%%%%%%%%%%%%%%%%%%%%%
\subsection{Numerical PDE Workflow}
\label{subsec:numerical-pde-workflow}
%%%%%%%%%%%%%%%%%%%%%%%%%%%%%%%%%%%%%%%%%%%%%%%%%%%%%%%%%%%%

A practical workflow is:

\begin{enumerate}
    \item identify the PDE and its mathematical type;
    \item define the computational domain;
    \item specify initial conditions when time dependent;
    \item specify all boundary conditions;
    \item identify conservation laws and invariants;
    \item identify expected regularity and possible discontinuities;
    \item choose finite difference, finite volume, finite element, spectral,
          or another discretization;
    \item construct the spatial mesh;
    \item choose the spatial approximation order;
    \item choose a time integrator if needed;
    \item derive stability restrictions;
    \item assemble the discrete algebraic system;
    \item exploit sparsity and matrix structure;
    \item select linear and nonlinear solvers;
    \item choose preconditioners;
    \item compute the numerical solution;
    \item monitor solver residuals;
    \item verify boundary conditions;
    \item check conservation or physical constraints;
    \item refine the mesh;
    \item refine the time step;
    \item estimate convergence rates;
    \item compare with exact or manufactured solutions;
    \item validate against physical data when applicable.
\end{enumerate}

%%%%%%%%%%%%%%%%%%%%%%%%%%%%%%%%%%%%%%%%%%%%%%%%%%%%%%%%%%%%
\subsection{Choosing a PDE Discretization}
\label{subsec:choosing-pde-discretization}
%%%%%%%%%%%%%%%%%%%%%%%%%%%%%%%%%%%%%%%%%%%%%%%%%%%%%%%%%%%%

For simple rectangular domains and smooth problems, consider:

\[
\boxed{
\text{finite differences}.
}
\]

For conservation laws and fluid transport, consider:

\[
\boxed{
\text{finite volumes}.
}
\]

For complex geometry and variational elliptic problems, consider:

\[
\boxed{
\text{finite elements}.
}
\]

For smooth periodic or simple-domain problems requiring very high
accuracy, consider:

\[
\boxed{
\text{spectral methods}.
}
\]

%%%%%%%%%%%%%%%%%%%%%%%%%%%%%%%%%%%%%%%%%%%%%%%%%%%%%%%%%%%%
\subsection{Choosing a Time Integrator}
\label{subsec:choosing-pde-time-integrator}
%%%%%%%%%%%%%%%%%%%%%%%%%%%%%%%%%%%%%%%%%%%%%%%%%%%%%%%%%%%%

For nonstiff hyperbolic systems, explicit methods are common.

For diffusive systems, explicit stability restrictions may force

\[
\boxed{
\Delta t=O(h^2).
}
\]

Implicit methods may permit much larger stable steps.

For method-of-lines systems, time-integration choices should follow the
ODE stability properties of the semi-discrete system.

%%%%%%%%%%%%%%%%%%%%%%%%%%%%%%%%%%%%%%%%%%%%%%%%%%%%%%%%%%%%
\subsection{Common Errors}
\label{subsec:numerical-pde-common-errors}
%%%%%%%%%%%%%%%%%%%%%%%%%%%%%%%%%%%%%%%%%%%%%%%%%%%%%%%%%%%%

\subsubsection{Using the Wrong PDE Classification}

Elliptic, parabolic, and hyperbolic equations have fundamentally different
mathematical behavior.

A method appropriate for diffusion may be inappropriate for wave
propagation.

%%%%%%%%%%%%%%%%%%%%%%%%%%%%%%%%%%%%%%%%%%%%%%%%%%%%%%%%%%%%
\subsubsection{Ignoring Boundary Conditions}

A PDE without correct boundary data may be incomplete or solve a different
physical problem.

Boundary discretization must be incorporated into the numerical system.

%%%%%%%%%%%%%%%%%%%%%%%%%%%%%%%%%%%%%%%%%%%%%%%%%%%%%%%%%%%%
\subsubsection{Using Centered Advection Without Stability Analysis}

Forward-time centered-space discretization of pure advection is unstable.

Higher formal symmetry does not imply stability.

%%%%%%%%%%%%%%%%%%%%%%%%%%%%%%%%%%%%%%%%%%%%%%%%%%%%%%%%%%%%
\subsubsection{Using an Explicit Diffusion Step That Violates the Stability Limit}

For one-dimensional FTCS heat flow,

\[
\boxed{
\frac{
\alpha\Delta t
}{
h^2
}
\leq
\frac12
}
\]

is required.

Violating this condition can create oscillatory growth.

%%%%%%%%%%%%%%%%%%%%%%%%%%%%%%%%%%%%%%%%%%%%%%%%%%%%%%%%%%%%
\subsubsection{Confusing Unconditional Stability With Unlimited Accuracy}

Backward Euler can be stable with large time steps.

That does not mean large time steps are accurate.

%%%%%%%%%%%%%%%%%%%%%%%%%%%%%%%%%%%%%%%%%%%%%%%%%%%%%%%%%%%%
\subsubsection{Ignoring the CFL Condition}

For explicit hyperbolic schemes, numerical information must propagate
quickly enough relative to the physical characteristics.

A violated CFL restriction can destroy stability.

%%%%%%%%%%%%%%%%%%%%%%%%%%%%%%%%%%%%%%%%%%%%%%%%%%%%%%%%%%%%
\subsubsection{Using the Wrong Upwind Direction}

For

\[
c>0,
\]

use left-biased information.

For

\[
c<0,
\]

use right-biased information.

Using the opposite direction produces a downwind scheme with poor
stability properties.

%%%%%%%%%%%%%%%%%%%%%%%%%%%%%%%%%%%%%%%%%%%%%%%%%%%%%%%%%%%%
\subsubsection{Interpreting Numerical Diffusion as Physical Diffusion}

A smeared numerical shock may be caused by the algorithm rather than the
physical PDE.

Grid refinement and method comparison can help separate the two.

%%%%%%%%%%%%%%%%%%%%%%%%%%%%%%%%%%%%%%%%%%%%%%%%%%%%%%%%%%%%
\subsubsection{Interpreting Numerical Oscillations as Physical Waves}

High-order dispersive schemes may generate oscillations near steep
gradients or discontinuities.

These may be numerical artifacts.

%%%%%%%%%%%%%%%%%%%%%%%%%%%%%%%%%%%%%%%%%%%%%%%%%%%%%%%%%%%%
\subsubsection{Ignoring Conservation}

For conservation laws, a nonconservative discretization can produce
incorrect shock speeds.

Finite-volume or conservative-difference formulations are often essential.

%%%%%%%%%%%%%%%%%%%%%%%%%%%%%%%%%%%%%%%%%%%%%%%%%%%%%%%%%%%%
\subsubsection{Ignoring Matrix Sparsity}

PDE discretization matrices are often sparse.

Treating them as dense can make a tractable problem computationally
impossible.

%%%%%%%%%%%%%%%%%%%%%%%%%%%%%%%%%%%%%%%%%%%%%%%%%%%%%%%%%%%%
\subsubsection{Using an Unpreconditioned Krylov Method on a Difficult PDE Matrix}

A Krylov solver may technically converge but require an impractical number
of iterations.

Preconditioning is often essential.

%%%%%%%%%%%%%%%%%%%%%%%%%%%%%%%%%%%%%%%%%%%%%%%%%%%%%%%%%%%%
\subsubsection{Refining Space Without Refining Time}

A spatial convergence study can be contaminated by time-discretization
error.

Both sources of error must be controlled.

%%%%%%%%%%%%%%%%%%%%%%%%%%%%%%%%%%%%%%%%%%%%%%%%%%%%%%%%%%%%
\subsubsection{Refining Time Without Refining Space}

Likewise, time error may become negligible while spatial error dominates.

The observed convergence may then stop improving.

%%%%%%%%%%%%%%%%%%%%%%%%%%%%%%%%%%%%%%%%%%%%%%%%%%%%%%%%%%%%
\subsubsection{Ignoring Boundary Accuracy}

A second-order interior stencil combined with a first-order boundary
formula may reduce the global convergence rate.

%%%%%%%%%%%%%%%%%%%%%%%%%%%%%%%%%%%%%%%%%%%%%%%%%%%%%%%%%%%%
\subsubsection{Assuming a Fine Mesh Automatically Gives a Correct Solution}

A fine mesh does not repair:

\[
\boxed{
\text{incorrect equations},
}
\]

\[
\boxed{
\text{incorrect boundary conditions},
}
\]

or

\[
\boxed{
\text{unstable discretization}.
}
\]

Verification requires more than resolution.

%%%%%%%%%%%%%%%%%%%%%%%%%%%%%%%%%%%%%%%%%%%%%%%%%%%%%%%%%%%%
\subsubsection{Using High-Order Spectral Methods Across Discontinuities}

Global smooth basis functions can produce Gibbs oscillations near jumps.

Shock-capturing or local methods may be more appropriate.

%%%%%%%%%%%%%%%%%%%%%%%%%%%%%%%%%%%%%%%%%%%%%%%%%%%%%%%%%%%%
\subsubsection{Ignoring Nonlinear Solver Convergence}

For nonlinear PDEs, discretization may be accurate while Newton iteration
has not converged.

Both discretization error and algebraic solver error must be monitored.

%%%%%%%%%%%%%%%%%%%%%%%%%%%%%%%%%%%%%%%%%%%%%%%%%%%%%%%%%%%%
\subsection{Applications}
\label{subsec:numerical-pde-applications}
%%%%%%%%%%%%%%%%%%%%%%%%%%%%%%%%%%%%%%%%%%%%%%%%%%%%%%%%%%%%

\begin{applicationbox}

\textbf{Heat Transfer.}
Diffusion and heat equations model temperature evolution in solids and
fluids.

\medskip

\textbf{Fluid Dynamics.}
Finite-volume and finite-element methods solve conservation laws and the
Navier--Stokes equations.

\medskip

\textbf{Elasticity.}
Finite elements are widely used for deformation, stress, and structural
mechanics.

\medskip

\textbf{Electromagnetics.}
Maxwell equations are solved numerically for antennas, waveguides, fields,
and electromagnetic devices.

\medskip

\textbf{Wave Propagation.}
Acoustic, seismic, and optical waves require stable low-dispersion
numerical methods.

\medskip

\textbf{Finance.}
Option-pricing equations such as Black--Scholes can be solved using
finite-difference and finite-element methods.

\medskip

\textbf{Biology.}
Reaction-diffusion equations model spatial population, chemical, and
pattern-forming systems.

\medskip

\textbf{Environmental Modeling.}
Transport-diffusion equations describe pollutants, heat, ocean processes,
and atmospheric transport.

\end{applicationbox}

%%%%%%%%%%%%%%%%%%%%%%%%%%%%%%%%%%%%%%%%%%%%%%%%%%%%%%%%%%%%
\subsection{Exercises}
\label{subsec:numerical-pde-exercises}
%%%%%%%%%%%%%%%%%%%%%%%%%%%%%%%%%%%%%%%%%%%%%%%%%%%%%%%%%%%%

\begin{exercise}[1]
Define a computational grid.
\end{exercise}

\begin{exercise}[1]
State the centered first-difference formula.
\end{exercise}

\begin{exercise}[1]
State the centered second-difference formula.
\end{exercise}

\begin{exercise}[1]
Define a finite-difference stencil.
\end{exercise}

\begin{exercise}[1]
Define a CFL condition conceptually.
\end{exercise}

\begin{exercise}[1]
Define a finite-volume cell average.
\end{exercise}

\begin{exercise}[1]
Define a weak formulation conceptually.
\end{exercise}

\begin{exercise}[1]
Define a finite-element basis function.
\end{exercise}

\begin{exercise}[1]
Define a stiffness matrix.
\end{exercise}

\begin{exercise}[1]
Define a spectral method.
\end{exercise}

\begin{exercise}[2]
Derive the centered second derivative using Taylor expansions.
\end{exercise}

\begin{exercise}[2]
Construct the three-point stencil for

\[
-u''=f.
\]
\end{exercise}

\begin{exercise}[3]
Build the finite-difference matrix for a one-dimensional Poisson problem
with four interior nodes.
\end{exercise}

\begin{exercise}[3]
Incorporate nonzero Dirichlet boundary values into the right-hand side.
\end{exercise}

\begin{exercise}[3]
Approximate a Neumann boundary condition using a one-sided difference.
\end{exercise}

\begin{exercise}[3]
Derive the five-point Laplacian.
\end{exercise}

\begin{exercise}[3]
Construct the two-dimensional Poisson equation at one interior node.
\end{exercise}

\begin{exercise}[3]
Explain why the resulting matrix is sparse.
\end{exercise}

\begin{exercise}[3]
Write the two-dimensional Laplacian using Kronecker products.
\end{exercise}

\begin{exercise}[2]
Derive the FTCS heat equation scheme.
\end{exercise}

\begin{exercise}[2]
Compute the diffusion number

\[
r
=
\frac{
\alpha\Delta t
}{
h^2
}.
\]
\end{exercise}

\begin{exercise}[3]
Determine whether a specified FTCS heat step is stable.
\end{exercise}

\begin{exercise}[3]
Find the maximum stable \(\Delta t\) for given \(h\) and \(\alpha\).
\end{exercise}

\begin{exercise}[3]
Derive backward Euler for the heat equation.
\end{exercise}

\begin{exercise}[3]
Write the resulting matrix system.
\end{exercise}

\begin{exercise}[3]
Derive Crank--Nicolson for the heat equation.
\end{exercise}

\begin{exercise}[3]
Compare the formal accuracy of FTCS, backward Euler, and
Crank--Nicolson.
\end{exercise}

\begin{exercise}[3]
Perform a von Neumann analysis of FTCS heat flow.
\end{exercise}

\begin{exercise}[2]
Define the Courant number for linear advection.
\end{exercise}

\begin{exercise}[3]
Derive the first-order upwind scheme for \(c>0\).
\end{exercise}

\begin{exercise}[3]
Derive the corresponding upwind scheme for \(c<0\).
\end{exercise}

\begin{exercise}[3]
Check the CFL condition for a specified advection problem.
\end{exercise}

\begin{exercise}[3]
Perform one upwind time step from given grid values.
\end{exercise}

\begin{exercise}[3]
Explain numerical diffusion in the upwind method.
\end{exercise}

\begin{exercise}[3]
Explain why centered FTCS is unstable for pure linear advection.
\end{exercise}

\begin{exercise}[3]
Apply one Lax--Friedrichs step.
\end{exercise}

\begin{exercise}[3]
Apply one Lax--Wendroff step.
\end{exercise}

\begin{exercise}[3]
Explain the difference between numerical diffusion and numerical
dispersion.
\end{exercise}

\begin{exercise}[2]
Derive the centered wave-equation scheme.
\end{exercise}

\begin{exercise}[3]
Determine its CFL stability condition.
\end{exercise}

\begin{exercise}[3]
Construct the first time level using initial displacement and velocity.
\end{exercise}

\begin{exercise}[3]
Perform one wave-equation update.
\end{exercise}

\begin{exercise}[2]
Integrate a conservation law over a finite-volume cell.
\end{exercise}

\begin{exercise}[3]
Derive the semi-discrete finite-volume update.
\end{exercise}

\begin{exercise}[3]
Explain why interface flux cancellation produces conservation.
\end{exercise}

\begin{exercise}[3]
Construct a consistent numerical flux for linear advection.
\end{exercise}

\begin{exercise}[3]
Explain the role of a Riemann problem in Godunov's method.
\end{exercise}

\begin{exercise}[3]
Compute discrete total variation for a given grid function.
\end{exercise}

\begin{exercise}[3]
Explain why TVD behavior is useful near shocks.
\end{exercise}

\begin{exercise}[2]
Derive the weak form of one-dimensional Poisson's equation.
\end{exercise}

\begin{exercise}[3]
Construct piecewise linear hat functions on a three-element mesh.
\end{exercise}

\begin{exercise}[3]
Compute a local finite-element stiffness matrix.
\end{exercise}

\begin{exercise}[3]
Assemble two local stiffness matrices into a global matrix.
\end{exercise}

\begin{exercise}[3]
Construct a finite-element mass matrix for linear basis functions.
\end{exercise}

\begin{exercise}[3]
Explain the difference between trial and test functions.
\end{exercise}

\begin{exercise}[3]
Explain the Galerkin condition.
\end{exercise}

\begin{exercise}[3]
Compare structured and unstructured meshes.
\end{exercise}

\begin{exercise}[3]
Explain \(h\)-refinement and \(p\)-refinement.
\end{exercise}

\begin{exercise}[2]
Write a Fourier spectral representation of a periodic function.
\end{exercise}

\begin{exercise}[3]
Show how differentiation acts on Fourier coefficients.
\end{exercise}

\begin{exercise}[3]
Explain spectral convergence.
\end{exercise}

\begin{exercise}[3]
Explain the Gibbs phenomenon.
\end{exercise}

\begin{exercise}[3]
Write a pseudospectral differentiation equation

\[
U_x\approx DU.
\]
\end{exercise}

\begin{exercise}[2]
Apply the method of lines to the heat equation.
\end{exercise}

\begin{exercise}[3]
Explain why the resulting ODE system becomes stiff as \(h\to0\).
\end{exercise}

\begin{exercise}[3]
Write a first-order operator splitting for a reaction-diffusion equation.
\end{exercise}

\begin{exercise}[3]
Write a Strang splitting sequence.
\end{exercise}

\begin{exercise}[3]
Construct a Newton linearization for a nonlinear discrete PDE system.
\end{exercise}

\begin{exercise}[3]
Explain Picard iteration.
\end{exercise}

\begin{exercise}[3]
Explain why Krylov methods are useful for PDE systems.
\end{exercise}

\begin{exercise}[3]
Explain the basic idea of multigrid.
\end{exercise}

\begin{exercise}[3]
Describe one two-grid correction cycle.
\end{exercise}

\begin{exercise}[3]
Explain the role of preconditioning in PDE linear solves.
\end{exercise}

\begin{exercise}[3]
Perform a mesh-refinement convergence study from supplied errors.
\end{exercise}

\begin{exercise}[3]
Estimate spatial convergence order.
\end{exercise}

\begin{exercise}[3]
Explain why both \(h\) and \(\Delta t\) must be refined in a
space-time convergence study.
\end{exercise}

\begin{exercise}[3]
Construct a manufactured solution for a Poisson problem.
\end{exercise}

\begin{exercise}[3]
Explain verification versus validation for a PDE simulation.
\end{exercise}

\begin{exercise}[4]
Derive the five-point Laplacian truncation error.
\end{exercise}

\begin{exercise}[4]
Prove positive definiteness of the one-dimensional Dirichlet discrete
Laplacian.
\end{exercise}

\begin{exercise}[4]
Study discrete maximum principles for elliptic finite differences.
\end{exercise}

\begin{exercise}[4]
Derive the eigenvalues of the one-dimensional discrete Laplacian.
\end{exercise}

\begin{exercise}[4]
Use those eigenvalues to analyze diffusion stiffness.
\end{exercise}

\begin{exercise}[4]
Perform a full von Neumann analysis of Crank--Nicolson.
\end{exercise}

\begin{exercise}[4]
Derive the instability of centered FTCS advection.
\end{exercise}

\begin{exercise}[4]
Perform modified-equation analysis for first-order upwind advection.
\end{exercise}

\begin{exercise}[4]
Study Lax--Wendroff dispersion.
\end{exercise}

\begin{exercise}[4]
Derive the CFL condition from numerical and physical domains of dependence.
\end{exercise}

\begin{exercise}[4]
Study weak solutions of scalar conservation laws.
\end{exercise}

\begin{exercise}[4]
Study entropy conditions and entropy-stable numerical methods.
\end{exercise}

\begin{exercise}[4]
Derive Godunov fluxes for scalar convex conservation laws.
\end{exercise}

\begin{exercise}[4]
Study MUSCL-type reconstruction.
\end{exercise}

\begin{exercise}[4]
Study high-resolution flux limiters.
\end{exercise}

\begin{exercise}[4]
Derive the weak finite-element formulation of Poisson's equation in two
dimensions.
\end{exercise}

\begin{exercise}[4]
Study the Lax--Milgram theorem as the foundation of elliptic finite
elements.
\end{exercise}

\begin{exercise}[4]
Study C\'ea's lemma.
\end{exercise}

\begin{exercise}[4]
Derive finite-element interpolation error estimates.
\end{exercise}

\begin{exercise}[4]
Study adaptive finite-element error estimation.
\end{exercise}

\begin{exercise}[4]
Study Fourier spectral convergence for analytic functions.
\end{exercise}

\begin{exercise}[4]
Study Chebyshev collocation methods.
\end{exercise}

\begin{exercise}[4]
Study multidimensional spectral methods.
\end{exercise}

\begin{exercise}[4]
Study algebraic multigrid.
\end{exercise}

\begin{exercise}[4]
Study domain-decomposition preconditioners.
\end{exercise}

\begin{exercise}[4]
Study stabilized finite elements for advection-dominated problems.
\end{exercise}

\begin{exercise}[4]
Study discontinuous Galerkin methods.
\end{exercise}

%%%%%%%%%%%%%%%%%%%%%%%%%%%%%%%%%%%%%%%%%%%%%%%%%%%%%%%%%%%%
\subsection{Conceptual Summary}
\label{subsec:numerical-pde-summary}
%%%%%%%%%%%%%%%%%%%%%%%%%%%%%%%%%%%%%%%%%%%%%%%%%%%%%%%%%%%%

Numerical PDE methods replace continuous derivatives and function spaces
with finite-dimensional approximations.

Finite differences use formulas such as

\[
\boxed{
u_{xx}(x_i)
\approx
\frac{
u_{i+1}-2u_i+u_{i-1}
}{
h^2
}.
}
\]

For two-dimensional diffusion,

\[
\boxed{
\Delta u_{i,j}
\approx
\frac{
u_{i+1,j}
+
u_{i-1,j}
+
u_{i,j+1}
+
u_{i,j-1}
-
4u_{i,j}
}{
h^2
}.
}
\]

The FTCS heat scheme satisfies

\[
\boxed{
u_i^{n+1}
=
r u_{i-1}^n
+
(1-2r)u_i^n
+
r u_{i+1}^n,
}
\]

with

\[
\boxed{
r
=
\frac{
\alpha\Delta t
}{
h^2
}.
}
\]

For stability in one dimension,

\[
\boxed{
r\leq\frac12.
}
\]

For linear advection,

\[
\boxed{
u_t+cu_x=0,
}
\]

upwind schemes use information from the characteristic direction and obey
CFL restrictions.

Finite-volume methods emphasize

\[
\boxed{
\text{conservation},
}
\]

finite-element methods emphasize

\[
\boxed{
\text{weak formulations and local basis functions},
}
\]

and spectral methods emphasize

\[
\boxed{
\text{high-order global approximation}.
}
\]

Reliable PDE computation requires simultaneous attention to

\[
\boxed{
\text{consistency}
+
\text{stability}
+
\text{convergence}
+
\text{conservation}
+
\text{linear algebra}
+
\text{mesh resolution}.
}
\]

%%%%%%%%%%%%%%%%%%%%%%%%%%%%%%%%%%%%%%%%%%%%%%%%%%%%%%%%%%%%
\subsection{Section Connections}
\label{subsec:numerical-pde-connections}
%%%%%%%%%%%%%%%%%%%%%%%%%%%%%%%%%%%%%%%%%%%%%%%%%%%%%%%%%%%%

\chapterconnection{
Numerical PDEs extend the approximation and stability ideas of
Section~\ref{sec:numerical-analysis} and the time-integration framework of
Section~\ref{sec:numerical-odes} to spatially distributed systems.
Finite-difference, finite-volume, finite-element, spectral, and
method-of-lines approaches transform PDEs into large algebraic or ODE
systems whose solution relies heavily on
Section~\ref{sec:numerical-linear-algebra}. Nonlinear PDEs also require
the iterative methods of Section~\ref{sec:numerical-optimization}. The
next section, Monte Carlo Methods, develops an entirely different
computational strategy based on random sampling for integration,
probability estimation, uncertainty quantification, simulation, and
high-dimensional computation.
}

%%%%%%%%%%%%%%%%%%%%%%%%%%%%%%%%%%%%%%%%%%%%%%%%%%%%%%%%%%%%
% SECTION SOURCE TRACKING
%%%%%%%%%%%%%%%%%%%%%%%%%%%%%%%%%%%%%%%%%%%%%%%%%%%%%%%%%%%%

%------------------------------------------------------------
% 11.5 Numerical Partial Differential Equations
%------------------------------------------------------------
%
% Source basis:
%   - Standard finite-difference PDE methods.
%   - Standard finite-volume, finite-element, and spectral methods.
%   - Standard PDE stability and convergence analysis.
%
% Used for:
%   - PDE classification;
%   - computational grids;
%   - finite-difference formulas;
%   - stencils;
%   - one- and two-dimensional Poisson equations;
%   - discrete Laplacians;
%   - Dirichlet, Neumann, and Robin conditions;
%   - sparse PDE matrices;
%   - Kronecker-product structure;
%   - discrete maximum-principle preview;
%   - heat equation;
%   - FTCS;
%   - BTCS;
%   - Crank--Nicolson;
%   - diffusion-number stability restrictions;
%   - von Neumann stability analysis;
%   - advection equation;
%   - Courant number;
%   - upwind methods;
%   - numerical diffusion;
%   - numerical dispersion;
%   - Lax--Friedrichs;
%   - Lax--Wendroff;
%   - CFL condition;
%   - wave-equation discretization;
%   - conservation laws;
%   - finite-volume methods;
%   - numerical fluxes;
%   - Godunov methods preview;
%   - shocks and weak solutions;
%   - monotone and TVD schemes;
%   - flux-limiters preview;
%   - weak formulations;
%   - finite-element methods;
%   - Galerkin approximation;
%   - stiffness and mass matrices;
%   - finite-element assembly;
%   - structured and unstructured meshes;
%   - adaptive mesh refinement;
%   - h-, p-, and hp-refinement;
%   - spectral methods;
%   - Fourier and pseudospectral methods;
%   - method of lines;
%   - semi-discrete stiffness;
%   - operator splitting;
%   - Strang splitting;
%   - nonlinear PDEs;
%   - Newton and Picard linearization;
%   - Krylov solvers;
%   - multigrid;
%   - preconditioning;
%   - domain decomposition;
%   - consistency, stability, and convergence;
%   - Lax equivalence preview;
%   - refinement studies;
%   - manufactured solutions;
%   - boundary errors and ghost points;
%   - nonuniform and anisotropic meshes;
%   - boundary layers;
%   - Peclet-number preview;
%   - advection-diffusion;
%   - reaction-diffusion;
%   - applications and exercises.
%
% Notes:
%   - Monte Carlo Methods are developed next in Section 11.6.
%   - Scientific Computing follows in Section 11.7.
%   - Numerical ODEs appear in Section 11.4.
%   - PDE theory appears in Chapter 9.
%
%%%%%%%%%%%%%%%%%%%%%%%%%%%%%%%%%%%%%%%%%%%%%%%%%%%%%%%%%%%%